%%
%% MScheme Compiler Design: From Binding Optimization to Native Code
%%
\documentclass[12pt]{report}
\usepackage[margin=1in]{geometry}
\usepackage{longtable}
\usepackage{booktabs}
\usepackage{listings}
\usepackage{xcolor}
\usepackage{hyperref}
\usepackage{textcomp}
\usepackage{amssymb}

\lstdefinelanguage{Scheme}{
  morekeywords={define,lambda,if,cond,else,let,let*,letrec,begin,
    set!,quote,quasiquote,unquote,unquote-splicing,and,or,not,
    case,do,delay,force,map,for-each,apply,call-with-current-continuation,
    call/cc,require-modules,load,macro,unwind-protect,eval,
    current-environment},
  sensitive=true,
  morecomment=[l]{;},
  morestring=[b]",
}

\lstdefinelanguage{Modula3}{
  morekeywords={MODULE,INTERFACE,IMPORT,FROM,PROCEDURE,BEGIN,END,
    VAR,CONST,TYPE,RECORD,OBJECT,METHODS,OVERRIDES,RETURN,
    IF,THEN,ELSE,ELSIF,WHILE,DO,FOR,TO,BY,LOOP,EXIT,REPEAT,UNTIL,
    CASE,OF,WITH,TRY,EXCEPT,FINALLY,RAISE,RAISES,
    NEW,ARRAY,REF,REFANY,SET,BOOLEAN,INTEGER,REAL,LONGREAL,CHAR,TEXT,
    TRUE,FALSE,NIL,BRANDED,REVEAL,EVAL,NOT,AND,OR,IN,
    ISTYPE,NARROW,TYPECODE,ADDRESS},
  sensitive=true,
  morecomment=[n]{(*}{*)},
  morestring=[b]",
}

\lstdefinelanguage{C}{
  morekeywords={typedef,char,int,void,if,else,return,while,for,
    extern,static,struct,union,const,unsigned,signed,long,double},
  sensitive=true,
  morecomment=[l]{//},
  morecomment=[s]{/*}{*/},
  morestring=[b]",
}

\lstset{
  basicstyle=\ttfamily\small,
  breaklines=true,
  frame=single,
  backgroundcolor=\color{gray!5},
  keywordstyle=\bfseries\color{blue!70!black},
  commentstyle=\itshape\color{green!50!black},
  stringstyle=\color{red!60!black},
  showstringspaces=false,
  columns=flexible,
  tabsize=2,
  xleftmargin=2em,
  framexleftmargin=1.5em,
}

\title{\bfseries\sffamily\fontsize{24.88}{30}\selectfont
  MScheme Compiler Design\\[0.5em]
  \large From Binding Optimization to Native Code Generation}
\author{Mika Nystr\"om \\ {\tt mika@alum.mit.edu}}
\date{March 2026}

\begin{document}
\maketitle
\thispagestyle{empty}

\newpage
\tableofcontents
\newpage

%======================================================================
\chapter{Introduction}
%======================================================================

The MScheme interpreter contains an optimization called the
\emph{binding optimization}, described in detail in the companion
report \emph{MScheme Internals}.  When a closure is created, the
\texttt{SchemeClosure.Init} procedure walks the body and replaces
free variable symbols with \texttt{SchemeEnvironmentBinding.T}
objects that point directly to the environment cell where each
variable lives.  At evaluation time, these binding objects are
recognized by the evaluator and resolved in constant time, bypassing
the scope-chain traversal.

This document explores the question: can this optimization be
extended into a framework for compiling Scheme closures to native
Modula-3 (or C) code?

We begin by analyzing what work the evaluator performs at runtime
that could instead be resolved statically (\S\ref{ch:analysis}),
then review existing infrastructure in CM3 that supports this goal
(\S\ref{ch:infra}), including a prior incomplete attempt from 2009
(\S\ref{sec:prior}).  We then propose three concrete compilation
strategies (\S\ref{ch:strategies}), classify Scheme constructs by
their compilability (\S\ref{ch:boundary}), and sketch the runtime
support needed (\S\ref{ch:runtime}).

%======================================================================
\chapter{Analysis: What the Eval Loop Does That a Compiler Could Eliminate}
\label{ch:analysis}
%======================================================================

The core evaluator, \texttt{EvalInternal} in \texttt{Scheme.m3},
repeats five categories of work on every iteration of its main loop.
All of them are partially or fully resolvable at compile time.

\section{Type Dispatch}
\label{sec:typedispatch}

Every iteration begins with a five-way type test:

\begin{lstlisting}[language=Modula3]
IF x = NIL THEN
  RETURN NIL
ELSIF ISTYPE(x, Symbol) THEN
  RETURN env.lookup(x)
ELSIF ISTYPE(x, Binding) THEN
  RETURN NARROW(x, Binding).get()
ELSIF NOT ISTYPE(x, Pair) THEN
  RETURN x                     (* self-evaluating *)
ELSE
  (* x is a Pair: special form or application *)
END
\end{lstlisting}

A compiler walking the source AST knows at compile time which branch
each sub-expression falls into.  A number literal is always
self-evaluating.  A symbol is always a variable reference.  A list
starting with \texttt{if} is always the \texttt{if} special form.
The type dispatch can be fully eliminated.

\subsection{Syntactic Keywords are Reserved}
\label{sec:keywords-reserved}

The compiler assumes that the R4RS syntactic keywords---\texttt{and},
\texttt{begin}, \texttt{case}, \texttt{cond}, \texttt{define},
\texttt{delay}, \texttt{do}, \texttt{if}, \texttt{lambda},
\texttt{let}, \texttt{let*}, \texttt{letrec}, \texttt{or},
\texttt{quote}, \texttt{set!}---have their standard meanings and are
never shadowed by variable bindings.  This is mandated by
R4RS~\S2.1:\footnote{William Clinger and Jonathan Rees~(eds.),
\emph{Revised$^4$ Report on the Algorithmic Language Scheme},
ACM Lisp Pointers~IV(3), July--September 1991.}

\begin{quote}
``Certain identifiers are reserved for use as syntactic keywords
and may not be used as variables.''
\end{quote}

R4RS goes on to note that ``some implementations allow all
identifiers, including syntactic keywords, to be used as variables''
but calls this out as creating ``ambiguities in the language.''
The MScheme interpreter is such an implementation: it permits
redefining \texttt{and}, \texttt{or}, \texttt{not}, \texttt{let},
\texttt{let*}, \texttt{case}, and~\texttt{do} as ordinary
procedures, while \texttt{if}, \texttt{begin}, \texttt{lambda},
\texttt{set!}, and~\texttt{quote} are hardcoded and cannot be
overridden.

The compiler follows R4RS and treats all syntactic keywords as
reserved.  Programs that redefine any of these names will produce
different results when compiled than when interpreted.

\section{Special-Form Recognition}
\label{sec:sfrecog}

When \texttt{x} is a pair, the evaluator extracts the first element
and compares it against 14 pre-interned special-form symbols:

\begin{lstlisting}[language=Modula3]
IF    fn = SYMquote THEN ...
ELSIF fn = SYMbegin THEN ...
ELSIF fn = SYMdefine THEN ...
ELSIF fn = SYMsetB THEN ...
ELSIF fn = SYMif THEN ...
ELSIF fn = SYMcond THEN ...
ELSIF fn = SYMlambda THEN ...
ELSIF fn = SYMmacro THEN ...
...
\end{lstlisting}

These are pointer comparisons against interned atoms---fast, but
unnecessary.  A compiler recognizes \texttt{if} syntactically and
emits the corresponding control flow directly.  No chain of runtime
comparisons is needed.

\section{Operator Type Dispatch}
\label{sec:opdispatch}

After recognizing that a form is an application (not a special form),
the evaluator performs a \texttt{TYPECASE} on the evaluated operator:

\begin{lstlisting}[language=Modula3]
TYPECASE fn OF
  Closure(c) => (* env setup, tail-call opt *)
| Macro(m)   => x := m.expand(t, x, args)
| Procedure(p) => RETURN p.apply(...)
ELSE RAISE E("Not a procedure")
END
\end{lstlisting}

For calls to known functions---which are the common case in
well-structured code---the operator type is statically known.
A call to a known closure can be compiled as a direct procedure
call.  A call to a known primitive (e.g., \texttt{+}) can be
inlined entirely.

\section{Primitive Dispatch}
\label{sec:primdispatch}

The primitive system uses a \texttt{CASE} statement over
$\sim$170 enumerated values:

\begin{lstlisting}[language=Modula3]
CASE VAL(t.id, P) OF
  P.Plus  => RETURN PrimPlus(args)
| P.Minus => RETURN PrimMinus(args)
| P.Car   => RETURN PedanticFirst(x)
| ...
END
\end{lstlisting}

When the compiler knows that a call is to the primitive \texttt{+},
it can emit the addition directly, bypassing the dispatch table,
the arity check, and the argument-list unpacking.

\section{Argument List Construction}
\label{sec:arglist}

The evaluator calls \texttt{evalList} to evaluate all arguments and
build a cons-cell list, which is then passed to the procedure's
\texttt{apply} method.  The \texttt{apply1}/\texttt{apply2}
fast paths mitigate this for small arities, but a compiled call
needs no intermediate list at all---arguments are passed directly
as Modula-3 procedure parameters or in a fixed-size array.

\section{Summary}

\begin{center}
\begin{tabular}{lll}
\toprule
\textbf{Eval loop work} & \textbf{Per-iteration cost} & \textbf{Compiled} \\
\midrule
Type dispatch & 3--5 ISTYPE checks & Eliminated \\
Special-form recognition & up to 14 pointer comparisons & Eliminated \\
Operator TYPECASE & 1 runtime type test & Eliminated for known calls \\
Primitive CASE dispatch & 1 table lookup & Inlined \\
Argument list (evalList) & $n$ cons allocations & Direct passing \\
\bottomrule
\end{tabular}
\end{center}

The binding optimization already eliminates variable-lookup overhead.
A compiler extends this to eliminate \emph{all} of the above.

%======================================================================
\chapter{Existing Infrastructure}
\label{ch:infra}
%======================================================================

Several pieces of infrastructure in the CM3 ecosystem and in MScheme
itself support the construction of a compiler.

\section{The \texttt{SchemeProcedure.T} Interface}
\label{sec:proctype}

All callable objects in MScheme extend \texttt{SchemeProcedure.T}:

\begin{lstlisting}[language=Modula3]
TYPE T <: Public;
Public = OBJECT METHODS
    apply(interp : Scheme.T; args : Object) : Object
      RAISES { E };
    apply1(interp : Scheme.T; a1 : Object) : Object
      RAISES { E };
    apply2(interp : Scheme.T; a1, a2 : Object) : Object
      RAISES { E };
  END;
\end{lstlisting}

A compiled Scheme function can be a new subtype of
\texttt{SchemeProcedure.T} whose \texttt{apply} method is generated
native code.  Since the evaluator dispatches through this interface,
compiled and interpreted closures can call each other seamlessly---no
bridging code is needed.

\section{Dynamic Loading: \texttt{m3dl}}
\label{sec:m3dl}

The CM3 package \texttt{m3-libs/m3dl} provides Modula-3 bindings to
POSIX dynamic linking:

\begin{lstlisting}[language=Modula3]
(* DL.i3 *)
PROCEDURE Open(fileName : TEXT; flag : INTEGER) : ADDRESS
  RAISES { OpenFail };
PROCEDURE Sym(handle : ADDRESS; symbol : TEXT) : ADDRESS
  RAISES { SymbolError };
PROCEDURE Close(handle : ADDRESS) : INTEGER;
\end{lstlisting}

This wraps \texttt{dlopen}/\texttt{dlsym}/\texttt{dlclose} and
can load any shared library (\texttt{.so}/\texttt{.dylib}) at
runtime.  On Windows, equivalent functionality is available via
\texttt{WinBase.LoadLibrary} and \texttt{WinBase.GetProcAddress}.

For Modula-3 shared libraries, \texttt{RTLinker.AddUnit} can
integrate a dynamically loaded module into the running runtime,
including type registration and module initialization.

\section{Shared Library Building}

CM3 fully supports building position-independent shared libraries on
all major platforms.  The \texttt{-fPIC} flag is applied by default
to all generated code.  On Darwin:

\begin{lstlisting}
cc -dynamiclib -o libfoo.dylib foo.o
\end{lstlisting}

On Linux:

\begin{lstlisting}
cc -shared -o libfoo.so foo.o
\end{lstlisting}

This means both C and Modula-3 code can be compiled into loadable
shared libraries at runtime.

\section{Conservative Garbage Collection}
\label{sec:gc}

The Modula-3 garbage collector scans the C stack conservatively.
Local variables in C functions that hold traced references
(\texttt{REFANY}, i.e., \texttt{ADDRESS} in the C backend) are
found by the stack scanner without any special annotations.  This
means that \textbf{compiled C code is automatically GC-safe}---no
write barriers, root registration, or GC cooperation protocol is
needed.

This is the single most important enabler for a Scheme-to-C
compilation strategy.

\section{The Stub Generator as a Model}
\label{sec:sstubgen}

The \texttt{sstubgen} tool and the \texttt{scodegen.scm} library
already demonstrate the pattern of generating Modula-3 source code
from Scheme.  The stub generator parses Modula-3 interfaces using
\texttt{m3tk}, constructs Scheme data structures representing the
types and procedures, and then uses Scheme code to emit \texttt{.m3}
files that bridge between Scheme and Modula-3.

A Scheme compiler would use the same pattern: walk a Scheme AST
and emit source code (Modula-3 or C) implementing the computation.

\section{The 2009 Compiler Prototype}
\label{sec:prior}

An incomplete compiler prototype exists in the tree at
\texttt{sstubgen/example/src/}, dating from June 2009.  It consists
of three files that together sketch a complete design.

\subsection{\texttt{CompilerHooks.i3}: The Runtime Framework}

This interface defines the calling convention for compiled procedures:

\begin{lstlisting}[language=Modula3]
TYPE
  Frame = ARRAY OF REFANY;

  LocalProcedure = PROCEDURE (
    VAR args : ARRAY OF REFANY;
    g : Getter; s : Setter;
    c : Copier; a : Adapter) : LocalProcedureResult
    RAISES { E };
\end{lstlisting}

The key design elements:

\begin{description}
\item[\texttt{Frame}] Arguments and locals are stored in a flat
  \texttt{ARRAY OF REFANY}.  Slot~0 is reserved for a pointer to
  the enclosing frame (the ``up pointer''), slots $1 \ldots N$ hold
  arguments, and slots $N{+}1 \ldots N{+}M$ hold local variables
  introduced by \texttt{define}.

\item[\texttt{Getter/Setter}] Callbacks for accessing variables in
  enclosing scopes by \texttt{(up, pos)} coordinates, where
  \texttt{up} is the number of enclosing frames to traverse and
  \texttt{pos} is the slot index within that frame.

\item[\texttt{Copier}] Copies frames from a given depth upward,
  used when a closure captures variables from an enclosing scope.

\item[\texttt{Adapter}] Converts the flat frame representation
  back into a \texttt{SchemeEnvironment.T}, needed when compiled
  code calls \texttt{current-environment} or falls back to the
  interpreter.

\item[\texttt{LocalProcedureResult}] A record supporting three
  outcomes: a direct result, a tail call to another
  \texttt{LocalProcedure} (with up to 10 arguments inline), or
  a tail call with a heap-allocated argument array.
\end{description}

\subsection{\texttt{CompiledStuff.m3}: Hand-Compiled Example}

This file shows what a compiled version of
\texttt{(define (f n) (if (= 0 n) 1 (f (- n 1))))} would look
like, in two versions:

\paragraph{Version 1: Tail call via result.}
Returns a \texttt{Result} record with \texttt{tailCall.f := CP}
(a self-reference), letting an outer trampoline loop handle the
tail call.

\paragraph{Version 2: Self-tail-call optimized.}
When the compiler detects that the tail call is to the same function,
it emits a \texttt{LOOP} that reassigns the arguments and iterates:

\begin{lstlisting}[language=Modula3]
VAR args2 := args;
BEGIN
  LOOP
    IF TruthO(PrimEquals(..., Constants[0], args2[1])) THEN
      RETURN Result { result := Constants[1] }
    ELSE
      args2 := ARRAY OF ... { PrimMinus(..., args2[1], Constants[1]) }
    END
  END
END
\end{lstlisting}

This is the same tail-call-via-loop pattern used by the interpreter's
\texttt{EvalInternal}, but applied at compile time to the specific
function body.

\subsection{\texttt{compiler.scm}: The Scheme-Side Compiler}

This file implements (incompletely) the analysis and code generation
phases:

\begin{itemize}
\item \textbf{Closure introspection}: Uses \texttt{modula-type-op}
  to access the \texttt{params}, \texttt{body}, and \texttt{env}
  fields of a live \texttt{SchemeClosure.T} object.

\item \textbf{Compile-time environments}: Maintains a list of
  frames, where each frame is a list of symbols.
  \texttt{search-var} resolves a symbol to \texttt{(frame-index
  . slot-index)} coordinates.

\item \textbf{Macro expansion}: Detects macros at compile time and
  expands them before analysis.

\item \textbf{Code generation}: The function
  \texttt{make-lambda-wrapper} emits Modula-3 procedure text with
  \texttt{Getter}/\texttt{Setter}/\texttt{Copier}/\texttt{Adapter}
  callbacks, argument unpacking, and arity checking.

\item \textbf{Bytecode sketch}: Functions \texttt{push-x} and
  \texttt{push-pair-x} emit a stack-machine instruction sequence
  (\texttt{push-constant}, \texttt{push-value}, \texttt{push-global},
  \texttt{push-primitive}, \texttt{pop}).
\end{itemize}

The prototype was not completed.  It handles \texttt{quote},
\texttt{begin}, \texttt{define}, \texttt{set!}, and \texttt{if}
(via desugaring to \texttt{cond}), but the \texttt{cond} and
\texttt{lambda} cases are empty stubs.


%======================================================================
\chapter{Proposed Compilation Strategies}
\label{ch:strategies}
%======================================================================

We identify three viable strategies, ranging from simplest to most
ambitious.

\section{Strategy A: Template-Based Compilation}
\label{sec:stratA}

This strategy extends the 2009 \texttt{CompilerHooks.i3} approach.
Rather than generating code per-closure, a fixed set of
pre-compiled Modula-3 template procedures are parameterized at
runtime by:

\begin{itemize}
\item A flat \texttt{ARRAY OF REFANY} frame for arguments and locals.
\item \texttt{Getter}/\texttt{Setter} callbacks for lexical variable
  access.
\item A constants array.
\item A bytecode or instruction sequence describing the computation.
\end{itemize}

The ``compilation'' step produces a compact representation (essentially
bytecode) that the template procedures interpret---but at far lower
overhead than the full eval loop, because type dispatch, special-form
recognition, and operator dispatch are already resolved.

\paragraph{Advantages.}
No external compiler needed at runtime.  No dynamic loading.
Can be implemented entirely within the existing MScheme build.
Incremental: can coexist with the interpreter.

\paragraph{Disadvantages.}
Still interpreted (at a lower level).  Limited optimization
opportunities.  The bytecode interpreter adds its own dispatch
overhead.

\section{Strategy B: Emit Modula-3 Source (Ahead-of-Time)}
\label{sec:stratB}

Use the pattern already established by \texttt{sstubgen} and
\texttt{scodegen.scm}: walk the Scheme AST, emit \texttt{.i3}
and \texttt{.m3} files defining \texttt{SchemeProcedure.T}
subtypes, and compile them into the application with \texttt{cm3}.

The generated procedure would look like:

\begin{lstlisting}[language=Modula3]
TYPE Compiled_f = SchemeProcedure.T OBJECT
    bindings : REF ARRAY OF SchemeEnvironmentBinding.T;
    constants : REF ARRAY OF SchemeObject.T;
  OVERRIDES
    apply  := Apply_f;
    apply1 := Apply1_f;
  END;

PROCEDURE Apply1_f(self : Compiled_f;
                   interp : Scheme.T;
                   n : SchemeObject.T) : SchemeObject.T
  RAISES { Scheme.E } =
VAR n_val := n;
BEGIN
  LOOP
    IF SchemeBoolean.TruthO(
         SchemeLongReal.NumEq(self.constants[0], n_val)) THEN
      RETURN self.constants[1]
    ELSE
      n_val := SchemeLongReal.Sub(n_val, self.constants[1])
    END
  END
END Apply1_f;
\end{lstlisting}

\paragraph{Advantages.}
Full native code.  The CM3 optimizer can further improve the output.
Type-safe: generated code is checked by the Modula-3 compiler.
Can leverage M3's exception model for \texttt{unwind-protect}.

\paragraph{Disadvantages.}
Ahead-of-time only---requires knowing which functions to compile
before the application is built.  Suitable for library code and
known hot paths, not for dynamically-defined functions.

\section{Strategy C: Emit C, Compile and Load Dynamically (JIT)}
\label{sec:stratC}

This is the most ambitious strategy.  At runtime, when a closure is
created (or when it becomes ``hot''):

\begin{enumerate}
\item Walk the closure's body AST, as \texttt{DoBindings} already does.
\item Emit a C source file implementing the function body.
\item Invoke the system C compiler: \texttt{cc -shared -O2 -o /tmp/f.dylib /tmp/f.c}
\item Load the shared library via \texttt{DL.Open}.
\item Look up the compiled function via \texttt{DL.Sym}.
\item Wrap the function pointer in a \texttt{CompiledProcedure}
  object (a subtype of \texttt{SchemeProcedure.T}).
\item Replace the interpreted closure with the compiled version.
\end{enumerate}

\subsection{The C Calling Convention}

A compiled C function would have the following signature:

\begin{lstlisting}[language=C]
typedef char* SchemeObject;  /* = REFANY = ADDRESS */

SchemeObject compiled_f(
    SchemeObject interp,      /* Scheme.T */
    SchemeObject *args,       /* evaluated argument array */
    int nargs,
    SchemeObject *free_vars,  /* pre-resolved bindings */
    SchemeObject *constants   /* literal pool */
);
\end{lstlisting}

The C code would call back into Modula-3 helper routines through
thin \texttt{<*EXTERNAL*>} wrappers for operations that cannot be
inlined:

\begin{lstlisting}[language=C]
/* Declared in a helper M3 module, callable from C */
extern SchemeObject MScm_Truth(SchemeObject x);
extern SchemeObject MScm_NumAdd(SchemeObject a, SchemeObject b);
extern SchemeObject MScm_NumEq(SchemeObject a, SchemeObject b);
extern SchemeObject MScm_Car(SchemeObject pair);
extern SchemeObject MScm_Cdr(SchemeObject pair);
extern SchemeObject MScm_Cons(SchemeObject a, SchemeObject b,
                              SchemeObject interp);
extern SchemeObject MScm_Apply(SchemeObject interp,
                               SchemeObject proc,
                               SchemeObject args);
extern SchemeObject MScm_EnvLookup(SchemeObject env,
                                   SchemeObject sym);
\end{lstlisting}

GC safety is free (see \S\ref{sec:gc}): the conservative stack
scanner finds all traced references in C stack frames.

\subsection{The Modula-3 Wrapper}

\begin{lstlisting}[language=Modula3]
TYPE CompiledProcedure = SchemeProcedure.T OBJECT
    cFunc     : ADDRESS;  (* function pointer *)
    constants : REF ARRAY OF SchemeObject.T;
    freeVars  : REF ARRAY OF SchemeEnvironmentBinding.T;
  OVERRIDES
    apply  := CompiledApply;
    apply1 := CompiledApply1;
    apply2 := CompiledApply2;
  END;
\end{lstlisting}

The \texttt{CompiledApply} method evaluates arguments (if not
already evaluated), marshals them into a C-compatible array, calls
the C function pointer, and returns the result.  Exception handling
stays in Modula-3: the C functions return a sentinel (e.g.,
\texttt{NULL}) to signal an error, and the wrapper raises
\texttt{Scheme.E}.

\paragraph{Advantages.}
True JIT compilation.  Can compile hot functions discovered at
runtime.  The system C compiler provides serious optimization
(\texttt{-O2}).  Compiled and interpreted code interoperate
seamlessly via \texttt{SchemeProcedure.T}.

\paragraph{Disadvantages.}
Compilation latency (invoking \texttt{cc} takes tens of
milliseconds).  Requires a C compiler on the target system.
Platform-specific shared library conventions.  Error reporting
from compiled code is more complex.

\subsection{Tiered Execution}

Compilation latency suggests a tiered approach:

\begin{enumerate}
\item All closures start interpreted (as today).
\item A call counter or profiling hook identifies hot closures.
\item Hot closures are compiled in a background thread.
\item Once compiled, the closure object is replaced with the
  compiled version.
\end{enumerate}

The existing \texttt{SchemeClosure.T} could be extended with a
call counter field, or a separate profiling primitive could be
provided.

%======================================================================
\chapter{The Compilability Boundary}
\label{ch:boundary}
%======================================================================

Not all Scheme constructs can be compiled to static code.  The
binding optimization in \texttt{SchemeClosure.Init} already
identifies the boundary through its bail-out conditions.
A compiler extends these conditions into a full classification.

\section{Easily Compilable Constructs}

\begin{center}
\begin{longtable}{lp{9cm}}
\toprule
\textbf{Construct} & \textbf{Compilation strategy} \\
\midrule
\endhead
Self-evaluating literals &
  Emit a reference into the constants array. \\
\texttt{quote} &
  Emit a reference to the quoted datum in the constants array. \\
\texttt{if} &
  Emit \texttt{IF \ldots\ THEN \ldots\ ELSE \ldots\ END}. \\
\texttt{cond} &
  Emit cascaded \texttt{IF}/\texttt{ELSIF}. \\
\texttt{begin} &
  Emit sequential statements. \\
\texttt{lambda} &
  Generate a new \texttt{CompiledProcedure} or
  \texttt{SchemeClosure.T} (with the captured environment). \\
\texttt{set!} on a known binding &
  Emit \texttt{binding.setB(val)} or a direct array write. \\
Variable reference (parameter) &
  Emit \texttt{args[i]}. \\
Variable reference (free, pre-resolved) &
  Emit \texttt{freeVars[i].get()} or a direct quick-array read. \\
Known primitive call &
  Fetch the primitive via \texttt{freeVars[i].get()}, then call.
  The binding indirection preserves redefinition semantics
  (see \S\ref{sec:redefinition}). \\
Known closure call &
  Fetch the closure via \texttt{freeVars[i].get()}, then call
  \texttt{.apply()}.  Same binding indirection.  The TYPECASE
  dispatch is not eliminated, but the eval-loop overhead
  (type dispatch, special-form checks, argument-list construction)
  is. \\
Self-tail-call &
  Emit \texttt{LOOP} with argument reassignment. \\
\bottomrule
\end{longtable}
\end{center}

\section{Compilable with Caveats}

\begin{center}
\begin{longtable}{lp{9cm}}
\toprule
\textbf{Construct} & \textbf{Difficulty} \\
\midrule
\endhead
\texttt{define} in body &
  Compilable as a local variable assignment, but subsequent
  expressions cannot have their bindings pre-resolved (same
  limitation as the current \texttt{DoBindings}).  The compiler
  must fall back to interpreter-style lookup for later references
  to the defined name. \\
\texttt{unwind-protect} &
  Compilable using Modula-3's \texttt{TRY}/\texttt{FINALLY}/\texttt{EXCEPT}
  or C's \texttt{setjmp}/\texttt{longjmp} (less clean). \\
\texttt{set!} on unknown binding &
  Must fall back to \texttt{env.set()} method call. \\
Calls to unknown functions &
  Must use the full \texttt{SchemeProcedure.T.apply()} dispatch.
  Still faster than the interpreter because the arguments are
  already evaluated and the operator is already resolved. \\
Closures capturing mutable variables &
  Need a shared mutable cell.  The existing
  \texttt{SchemeEnvironmentBinding.T} with \texttt{get()}/\texttt{setB()}
  already provides this, but the direct quick-array optimization
  is lost when variables are captured across closure boundaries. \\
Variadic functions &
  Need runtime list construction for the rest parameter. \\
\bottomrule
\end{longtable}
\end{center}

\section{Must Remain Interpreted}

\begin{center}
\begin{longtable}{lp{9cm}}
\toprule
\textbf{Construct} & \textbf{Reason} \\
\midrule
\endhead
\texttt{eval} &
  By definition requires the interpreter.  Can be compiled as a
  call to \texttt{interp.eval()}, but the argument expression is
  unknown at compile time. \\
\texttt{current-environment} &
  A compiled function uses native variables, not a
  \texttt{SchemeEnvironment.Instance}.  Materializing an environment
  on demand requires the \texttt{Adapter} callback from
  \texttt{CompilerHooks.i3}---feasible but complex. \\
Top-level \texttt{define} &
  Adds bindings to the global environment at runtime.  The
  \texttt{define} itself compiles as \texttt{env.define(sym, val)},
  but the newly-defined name cannot be inlined by already-compiled
  code. \\
Reflective closure access &
  Code that introspects closure internals (e.g.,
  \texttt{get-proc-body}) expects the \texttt{params}/\texttt{body}
  fields of \texttt{SchemeClosure.T}.  A compiled closure does not
  have a meaningful \texttt{body} field. \\
\bottomrule
\end{longtable}
\end{center}

\section{The Key Invariant}

The binding optimization's bail-out conditions already identify the
compilability boundary:

\begin{itemize}
\item If the operator is a special form $\Rightarrow$ compile the
  corresponding control-flow pattern.
\item If the operator is a known procedure $\Rightarrow$ compile a
  call through the binding (see \S\ref{sec:redefinition}).
\item If the operator is a macro $\Rightarrow$ expand at compile time,
  then compile the expansion.
\item If the operator is unknown (not bound at closure-creation time)
  $\Rightarrow$ fall back to runtime dispatch.
\item If \texttt{define} appears $\Rightarrow$ stop optimizing
  subsequent expressions (or fall back to runtime for new bindings).
\end{itemize}

The compiler need not handle every case.  Any expression it cannot
compile can be delegated to the interpreter by emitting a call to
\texttt{interp.eval(quoted\_expr, env)}.  This provides a graceful
degradation path: the compiler handles what it can, and the
interpreter handles the rest.

%======================================================================
\chapter{Redefinition Semantics}
\label{ch:redef}
%======================================================================

A Scheme program may redefine any top-level symbol at any time.
When \texttt{foo} is redefined, all subsequent calls to \texttt{foo}
must dispatch to the new definition.  This is a fundamental semantic
requirement, and the compiler must preserve it.  Both compiled and
interpreted callers of \texttt{foo} must see the new value.

\section{How the Interpreter Handles Redefinition}

In the current interpreter, redefinition is always visible because
of how the binding optimization works.  A
\texttt{SchemeEnvironmentBinding.T} object points to the
\emph{cell} in the environment where the variable is stored, not
to the \emph{value} itself.  When \texttt{(set!\ foo new-proc)} or
\texttt{(define foo new-proc)} mutates the cell, every subsequent
call to \texttt{binding.get()} returns the new value.

For closures that were not subject to the binding optimization
(because the symbol was unbound at closure-creation time), the
evaluator falls back to \texttt{env.lookup(sym)}, which traverses
the live environment and also sees the new value.

In both cases, redefinition is immediately visible to all clients.

\section{Three Levels of Compilation Aggressiveness}
\label{sec:redefinition}

A compiler for Scheme must choose how aggressively it resolves
call targets.  There are three levels, with increasing performance
but decreasing compatibility with redefinition:

\begin{description}
\item[Level 1: Call through the Binding.]
  Compiled code fetches the current value of the callee from the
  \texttt{SchemeEnvironmentBinding.T} on every call:

\begin{lstlisting}[language=Modula3]
(* compiled call to foo *)
VAR proc := self.freeVars[3].get();
BEGIN
  RETURN NARROW(proc, SchemeProcedure.T).apply(interp, args)
END
\end{lstlisting}

  The Binding points to the environment cell.  If \texttt{foo} is
  redefined, the next call sees the new value.  This is semantically
  identical to the interpreter.  Both compiled and interpreted
  callers of \texttt{foo} see the redefinition immediately.

  The cost is one pointer dereference plus one indexed array read
  (for quick-array bindings) per call---negligible compared to the
  eval-loop overhead being eliminated.

\item[Level 2: Direct call.]
  The compiler resolves the callee at compile time and emits a
  direct procedure call:

\begin{lstlisting}[language=Modula3]
(* compiled direct call to foo's compiled body *)
RETURN Apply1_foo(self_foo, interp, arg1)
\end{lstlisting}

  This is faster (no indirection), but \textbf{broken under
  redefinition}: the old code is baked into the caller.  Redefining
  \texttt{foo} has no effect on this caller.

\item[Level 3: Inlining.]
  The compiler splices the callee's body into the caller.
  Redefinition is completely invisible---the old logic is
  permanently embedded.
\end{description}

\section{Design Decision: Level 1 Only}

This compiler adopts \textbf{Level 1 exclusively}.  All calls to
free variables go through \texttt{SchemeEnvironmentBinding.T}
objects.  No call targets are resolved beyond the binding cell.
No inlining across closure boundaries is performed.

This decision has several consequences:

\begin{enumerate}
\item \textbf{Redefinition always works.}  The compiler preserves
  the same semantics as the interpreter.  A user who redefines
  \texttt{foo} at the REPL sees the new definition take effect
  immediately, regardless of whether callers are compiled or
  interpreted.

\item \textbf{Compiled and interpreted code interoperate freely.}
  A compiled caller can call an interpreted callee (by fetching it
  via \texttt{.get()} and calling \texttt{.apply()}).  An
  interpreted caller can call a compiled callee (it is just a
  \texttt{SchemeProcedure.T} subtype).  Redefinition can replace
  a compiled procedure with an interpreted one or vice versa, and
  all clients adapt transparently.

\item \textbf{The performance benefit is still substantial.}
  The overhead eliminated by compilation is not the one pointer
  dereference per call (which Level~2 would remove).  It is the
  five-way type dispatch, the 14 special-form comparisons, the
  operator TYPECASE, the primitive CASE dispatch, and the
  argument-list cons-cell allocation---all of which are eliminated
  by Level~1 compilation.

\item \textbf{The \texttt{SchemeEnvironmentBinding.T} abstraction
  is reused.}  The same Binding objects that \texttt{DoBindings}
  already creates for the interpreter's optimization serve as the
  indirection cells for compiled code.  No new mechanism is needed.

\item \textbf{Level 2 and Level 3 remain available as future
  optimizations.}  If profiling reveals that a specific call site
  is a bottleneck \emph{and} the user can guarantee that the callee
  will not be redefined (e.g., via an explicit \texttt{(declare
  (constant foo))} annotation), the compiler could be extended to
  emit direct calls or inlined code for those specific sites.
  But this is strictly opt-in and is not part of the initial design.
\end{enumerate}

\section{Self-Tail-Calls: An Exception}

One case where the compiler \emph{can} safely bypass the binding
is self-tail-calls.  When a function calls itself in tail position,
the compiler knows the target is the same compiled body and can
emit a \texttt{LOOP} with argument reassignment.  This does not
violate redefinition semantics: if the function redefines itself
during execution (via \texttt{set!} on its own name), the
recursive call was already in progress and the new definition
takes effect on the \emph{next} external call, not the current
tail-recursive loop.  This matches the interpreter's behavior,
where \texttt{EvalInternal}'s \texttt{LOOP} evaluates
\texttt{c.body} (captured at closure-creation time), not a fresh
lookup of the function's name.

The same optimization applies to named \texttt{let} loops:
\texttt{(let loop ((i 0)) \ldots\ (loop (+ i 1)))} compiles as
a \texttt{LOOP} with variable reassignment, bypassing the
environment entirely.  Since the loop name is lexically scoped to
the \texttt{let} body, it is never visible in the environment and
cannot be redefined, so there is no semantic concern.

\paragraph{Tail position.}
The \texttt{has-self-tail-call?} predicate determines whether a
function body contains a self-call in tail position.  The following
forms propagate tail position to their sub-expressions:
\texttt{if} (consequent and alternative),
\texttt{cond} (last expression of each clause),
\texttt{case} (last expression of each clause),
\texttt{begin} (last expression),
\texttt{let}, \texttt{let*}, \texttt{letrec} (body),
\texttt{do} (exit expressions),
\texttt{and} (last argument),
and \texttt{or} (last argument).
The treatment of \texttt{and}/\texttt{or} as tail-propagating is
valid because R4RS reserves these as syntactic keywords
(see~\S\ref{sec:keywords-reserved}): they always have their
standard short-circuit semantics, so the last argument is genuinely
in tail position.  If they could be redefined as procedures, the
last argument would be an operand evaluated before the call, not
a tail expression.

\paragraph{Non-tail self-recursion.}
By contrast, when a function calls itself in \emph{non}-tail
position (e.g., \texttt{(* n (factorial (- n 1)))}), the call
goes through the binding as usual.  If \texttt{factorial} is
redefined between two recursive invocations, the new definition
\emph{is} visible.  The test suite exercises this: it redefines
callees between calls and verifies that compiled callers see the
new definitions.

%======================================================================
\chapter{Runtime Support}
\label{ch:runtime}
%======================================================================

\section{The \texttt{CompiledProcedure} Type}

The central new type is a subtype of \texttt{SchemeProcedure.T}:

\begin{lstlisting}[language=Modula3]
TYPE CompiledProcedure = SchemeProcedure.T OBJECT
    constants : REF ARRAY OF SchemeObject.T;
    freeVars  : REF ARRAY OF SchemeEnvironmentBinding.T;
  OVERRIDES
    apply  := CompiledApply;
    apply1 := CompiledApply1;
    apply2 := CompiledApply2;
  END;
\end{lstlisting}

For Strategy C (JIT), an additional \texttt{cFunc : ADDRESS} field
holds the dynamically-loaded function pointer.

The \texttt{constants} array holds all literal values referenced
by the function body: numbers, strings, quoted data, and interned
symbols.  These are allocated once and shared across all invocations.

The \texttt{freeVars} array holds \texttt{SchemeEnvironmentBinding.T}
objects for captured free variables.  These are the same binding
objects that \texttt{DoBindings} already creates---the compiler
simply stores them in an array rather than embedding them in the
AST.

\section{Interpreter Fallback}

Any compiled procedure that encounters a construct it cannot handle
natively falls back to the interpreter:

\begin{lstlisting}[language=Modula3]
(* In generated code: *)
result := interp.eval(quotedExpr, env)
\end{lstlisting}

For Strategy C, the C code calls a helper:

\begin{lstlisting}[language=C]
result = MScm_Eval(interp, quoted_expr, env);
\end{lstlisting}

This requires materializing a \texttt{SchemeEnvironment.T} from the
compiled frame.  The \texttt{Adapter} callback in
\texttt{CompilerHooks.i3} was designed for this: it converts the flat
frame array into a proper environment object that the interpreter
can traverse.

\section{Trampoline for Tail Calls}

For Strategy A and Strategy C, a trampoline loop is needed to handle
tail calls between compiled procedures without growing the stack:

\begin{lstlisting}[language=Modula3]
PROCEDURE Trampoline(proc : CompiledProcedure;
                     args : REF ARRAY OF SchemeObject.T;
                     interp : Scheme.T) : SchemeObject.T
  RAISES { Scheme.E } =
VAR result : LocalProcedureResult;
BEGIN
  result := proc.call(args);
  WHILE result.tailProc # NIL DO
    result := result.tailProc.call(result.tailArgs)
  END;
  RETURN result.result
END Trampoline;
\end{lstlisting}

Self-tail-calls (the most common case) are optimized to a loop
within the compiled function body itself, as shown in
\texttt{CompiledStuff.m3}.

\section{Exception Handling}

Scheme exceptions (\texttt{Scheme.E}) are Modula-3 exceptions.
In Strategies A and B, compiled code raises and catches them
natively.  In Strategy C (C code), exceptions require a protocol:

\begin{enumerate}
\item Helper functions return a sentinel value (e.g., a specific
  tagged pointer) to signal an exception.
\item The C function checks the return value and returns the
  sentinel to its caller.
\item The Modula-3 wrapper in \texttt{CompiledApply} checks the
  final result and raises \texttt{Scheme.E} if needed.
\end{enumerate}

Alternatively, \texttt{setjmp}/\texttt{longjmp} could be used for
non-local exits from C code, but this interacts poorly with the
Modula-3 exception model and is not recommended.

%======================================================================
\chapter{Recommended Approach}
\label{ch:recommendation}
%======================================================================

We recommend a phased implementation:

\paragraph{Phase 1: Strategy B (Ahead-of-Time).}
Implement a Scheme-to-Modula-3 code generator, driven from Scheme
(like \texttt{scodegen.scm}), that compiles known library functions
and hot paths at build time.  This requires:

\begin{itemize}
\item A compiler pass extending \texttt{DoBindings} to emit M3 source
  rather than transform the AST in place.
\item A \texttt{CompiledProcedure} type in a new \texttt{mscheme-compiler}
  package.
\item Quake build templates (analogous to \texttt{SchemeStubs} in
  \texttt{schemestubs.tmpl}) that invoke the compiler during the
  build.
\end{itemize}

This phase delivers immediate performance benefits for numerical
and list-processing code with no runtime complexity.

\paragraph{Phase 2: Strategy C (JIT).}
Add dynamic compilation for closures identified as hot at runtime.
This requires:

\begin{itemize}
\item A Scheme-to-C emitter (reusing the Phase 1 analysis).
\item A thin M3 helper library with \texttt{<*EXTERNAL*>} wrappers
  for the primitives.
\item Integration with \texttt{m3dl} for loading compiled code.
\item A call-counting or profiling mechanism to identify hot closures.
\item A compilation thread pool to avoid blocking the interpreter.
\end{itemize}

\paragraph{Phase 3: Optimization.}
Once the basic compilation pipeline works:

\begin{itemize}
\item Type specialization for numeric code (avoid boxing when the
  compiler can prove a variable is always a \texttt{LONGREAL}).
\item Inlining of small known callees.
\item Constant folding and dead-code elimination.
\item Escape analysis to stack-allocate environments that do not
  escape (extending the \texttt{assigned} flag mechanism).
\end{itemize}

%======================================================================
\chapter{Build System Integration}
\label{ch:build}
%======================================================================

The CM3 build system already supports a pattern for generating
Modula-3 source code at build time via Quake templates.  The
\texttt{SchemeStubs} mechanism in \texttt{sstubgen} uses this pattern
to generate M3 stubs from M3 interfaces.  The Scheme compiler can
use the same pattern in reverse: generating M3 code from Scheme
source.

\section{The \texttt{SchemeStubs} Pattern}

The existing stub-generation system uses a three-phase build
protocol, driven by Quake directives in the \texttt{m3makefile}:

\begin{description}
\item[Phase 1: \texttt{SchemeStubs("Foo")}]
  Runs the \texttt{sstubgen} tool on the interface \texttt{Foo.i3}.
  The tool parses the M3 interface using \texttt{m3tk}, walks the
  AST, and emits \texttt{FooSchemeStubs.i3} and
  \texttt{FooSchemeStubs.m3} into the build directory.  The Quake
  directive then calls \texttt{derived\_interface} and
  \texttt{derived\_implementation} to register the generated files
  for compilation by cm3.

\item[Phase 2: \texttt{ExportSchemeStubs("prog")}]
  Aggregates all per-interface stub modules generated in Phase~1
  into a package-level registration module
  \texttt{prog\_\_SCHEMEEXPORTS\_\_.m3}, which calls each stub
  module's \texttt{RegisterStubs()} in sequence.

\item[Phase 3: \texttt{importSchemeStubs()}]
  Scans the link context (\texttt{.M3IMPTAB}) for all packages
  that exported Scheme stubs, and generates a top-level aggregator
  \texttt{SchemeStubs.m3} that chains all packages' registrations.
  The executable's \texttt{Main} module calls
  \texttt{SchemeStubs.RegisterStubs()} at startup.
\end{description}

The key CM3 build-system primitives used are:
\begin{itemize}
\item \texttt{\_exec([...])}: runs an external tool during the Quake
  evaluation phase (before cm3 compiles anything).
\item \texttt{stale(target, source)}: skips re-generation if the
  output is newer than the input (incremental builds).
\item \texttt{derived\_interface(name, VISIBLE)} and
  \texttt{derived\_implementation(name)}: register generated source
  files for compilation by cm3.
\item \texttt{template("foo")}: loads \texttt{foo.tmpl} and makes
  its Quake procedures available.
\end{itemize}

\section{The \texttt{scheme()} Directive}

A new Quake template, \texttt{schemecompile.tmpl}, would define
the \texttt{scheme()} directive.  The three-phase protocol maps
directly onto the existing pattern:

\subsection{Phase 1: \texttt{scheme("mylib")}}

\begin{lstlisting}
readonly proc scheme(nam) is
    local src = path_of(nam & ".scm")
    local out = nam & "Compiled"

    if stale(out & ".i3", src)
        > "schemecompile.scm" in
            write ("(compile-scheme-file \""
                   & nam & "\")", CR)
        end
        _exec([_SCHEMECOMPILER,
               "-nointeractive",
               "-scm", _SCHEMECOMPILERSCM,
               "schemecompile.scm", _SEXIT])
    end

    compiled_intfs = compiled_intfs & " " & out
    derived_interface(out, VISIBLE)
    derived_implementation(out)
end
\end{lstlisting}

This runs the Scheme compiler tool on \texttt{mylib.scm} and emits
\texttt{mylibCompiled.i3} and \texttt{mylibCompiled.m3} into the
build directory.  The \texttt{stale} guard ensures the compiler is
only re-run when the Scheme source changes.

The compiler tool is itself an MScheme program---it reads
\texttt{.scm} files, walks the S-expression AST, and emits M3
source text using the same string-building pattern that
\texttt{sstubgen.scm} and \texttt{scodegen.scm} already use.

\subsection{Phase 2: \texttt{ExportSchemeCompiled("prog")}}

Aggregates all \texttt{scheme()} outputs from the current package
into a single registration module:

\begin{lstlisting}
readonly proc ExportSchemeCompiled(package) is
    local tgt = package & "__SCHEMECOMPILED__"
    (* generate tgt.i3 and tgt.m3 that call
       RegisterCompiled() on each compiled module *)
    ...
    derived_interface(tgt, VISIBLE)
    derived_implementation(tgt)
end
\end{lstlisting}

\subsection{Phase 3: \texttt{importSchemeCompiled()}}

Scans the link context for all packages that exported compiled
Scheme code and generates a top-level aggregator, exactly as
\texttt{importSchemeStubs()} does.

\section{Generated Code Structure}

For a Scheme source file \texttt{mylib.scm} containing:

\begin{lstlisting}[language=Scheme]
(define (add1 x) (+ x 1))

(define (apply-foo x) (foo x))

(define (factorial n)
  (if (= n 0)
      1
      (* n (factorial (- n 1)))))
\end{lstlisting}

The compiler would emit \texttt{mylibCompiled.i3}:

\begin{lstlisting}[language=Modula3]
INTERFACE mylibCompiled;
IMPORT SchemeEnvironment;
PROCEDURE RegisterCompiled(env : SchemeEnvironment.T);
END mylibCompiled.
\end{lstlisting}

And \texttt{mylibCompiled.m3}:

\begin{lstlisting}[language=Modula3]
MODULE mylibCompiled;
IMPORT Scheme, SchemeSymbol, SchemeProcedure,
       SchemeEnvironmentBinding, SchemeBoolean,
       SchemeLongReal, SchemeEnvironment;
FROM Scheme IMPORT E, Object;

(* --- Constants --- *)
VAR
  C0 : Object; (* 1, as SchemeLongReal *)
  C1 : Object; (* 0, as SchemeLongReal *)

(* === add1 === *)
TYPE Compiled_add1 = SchemeProcedure.T OBJECT
    k : REF ARRAY OF Object;           (* constants *)
    fv : REF ARRAY OF
             SchemeEnvironmentBinding.T; (* free vars *)
  OVERRIDES
    apply1 := Apply1_add1;
  END;

PROCEDURE Apply1_add1(self : Compiled_add1;
                      interp : Scheme.T;
                      x : Object) : Object
  RAISES { E } =
BEGIN
  (* (+ x 1) -- + is a known primitive *)
  RETURN SchemeLongReal.Add(x, self.k[0])
END Apply1_add1;

(* === apply-foo === *)
TYPE Compiled_apply_foo = SchemeProcedure.T OBJECT
    fv : REF ARRAY OF
             SchemeEnvironmentBinding.T;
  OVERRIDES
    apply1 := Apply1_apply_foo;
  END;

PROCEDURE Apply1_apply_foo(self : Compiled_apply_foo;
                           interp : Scheme.T;
                           x : Object) : Object
  RAISES { E } =
VAR proc := self.fv[0].get(); (* fetch foo via Binding *)
BEGIN
  RETURN NARROW(proc,
    SchemeProcedure.T).apply1(interp, x)
END Apply1_apply_foo;

(* === factorial === *)
TYPE Compiled_factorial = SchemeProcedure.T OBJECT
    k : REF ARRAY OF Object;
    fv : REF ARRAY OF
             SchemeEnvironmentBinding.T;
  OVERRIDES
    apply1 := Apply1_factorial;
  END;

PROCEDURE Apply1_factorial(self : Compiled_factorial;
                           interp : Scheme.T;
                           n : Object) : Object
  RAISES { E } =
BEGIN
  IF SchemeBoolean.TruthO(
       SchemeLongReal.NumEq(n, self.k[1])) THEN
    RETURN self.k[0]  (* 1 *)
  ELSE
    (* n * factorial(n - 1) *)
    VAR n_minus_1 := SchemeLongReal.Sub(n, self.k[0]);
        (* fetch factorial via Binding -- Level 1 *)
        fact_proc := self.fv[0].get();
        sub_result := NARROW(fact_proc,
          SchemeProcedure.T).apply1(interp, n_minus_1);
    BEGIN
      RETURN SchemeLongReal.Mul(n, sub_result)
    END
  END
END Apply1_factorial;

(* === Registration === *)
PROCEDURE RegisterCompiled(
    env : SchemeEnvironment.T) =
BEGIN
  C0 := SchemeLongReal.FromLR(1.0d0);
  C1 := SchemeLongReal.FromLR(0.0d0);

  WITH k = NEW(REF ARRAY OF Object, 1) DO
    k[0] := C0;
    env.define(SchemeSymbol.Symbol("add1"),
      NEW(Compiled_add1, k := k,
          name := "add1"))
  END;

  WITH fv = NEW(REF ARRAY OF
       SchemeEnvironmentBinding.T, 1) DO
    fv[0] := env.bind(SchemeSymbol.Symbol("foo"));
    env.define(SchemeSymbol.Symbol("apply-foo"),
      NEW(Compiled_apply_foo, fv := fv,
          name := "apply-foo"))
  END;

  WITH k = NEW(REF ARRAY OF Object, 2),
       fv = NEW(REF ARRAY OF
            SchemeEnvironmentBinding.T, 1) DO
    k[0] := C0;  k[1] := C1;
    env.define(SchemeSymbol.Symbol("factorial"),
      NEW(Compiled_factorial, k := k, fv := fv,
          name := "factorial"));
    (* Now that factorial is defined, bind fv[0]
       to its own Binding for self-tail-call *)
    fv[0] := env.bind(
               SchemeSymbol.Symbol("factorial"))
  END
END RegisterCompiled;

BEGIN END mylibCompiled.
\end{lstlisting}

\section{Key Design Points}

\begin{enumerate}
\item \textbf{The compiler tool is an MScheme program.}
  It reads \texttt{.scm} files, walks the AST (extending the
  analysis that \texttt{DoBindings} already performs), and emits
  Modula-3 source text.  This follows the established pattern:
  \texttt{sstubgen} is an MScheme program that reads M3 interfaces
  and emits M3 stubs; the Scheme compiler reads Scheme source and
  emits M3 compiled procedures.

\item \textbf{Binding capture happens at registration time.}
  The \texttt{RegisterCompiled} procedure calls \texttt{env.bind()}
  on the live global environment to obtain
  \texttt{SchemeEnvironmentBinding.T} objects.  This is identical
  to the moment when \texttt{SchemeClosure.Init} calls
  \texttt{DoBindings} for an interpreted closure---the same
  environment, the same Binding objects, the same semantics.

\item \textbf{Incremental builds.}  The \texttt{stale()} guard
  ensures the Scheme compiler is only re-run when the \texttt{.scm}
  source changes.  The generated \texttt{.m3} files are then
  recompiled by cm3 only if they changed (the standard CM3
  incremental build).

\item \textbf{Mixed compiled and interpreted definitions.}
  A single package can contain both \texttt{scheme("mylib")} (compiled)
  and \texttt{.scm} files loaded at runtime via \texttt{loadFile}
  (interpreted).  The compiled definitions are installed by
  \texttt{RegisterCompiled} at startup; interpreted definitions are
  loaded later.  Both coexist in the same global environment.

\item \textbf{Interpackage dependencies.}  The three-phase
  aggregation ensures that all compiled Scheme code from all
  transitively imported packages is registered before the
  application's \texttt{Main} module runs.  Module initialization
  order is handled by CM3's \texttt{RTLinker}, which topologically
  sorts the dependency graph.

\item \textbf{No new tools beyond MScheme itself.}  The compiler
  tool is built as an MScheme program using the existing
  \texttt{mscheme-interactive} infrastructure.  The Quake template
  invokes it via \texttt{\_exec}, just as \texttt{SchemeStubs}
  invokes \texttt{sstubgen}.
\end{enumerate}

\section{Example \texttt{m3makefile}}

A package using both Scheme compilation and M3 stub generation:

\begin{lstlisting}
import("sstubgen")
import("mscheme")
import("modula3scheme")
import("schemecompile")     (* new: loads the template *)

template("schemestubs")
template("schemecompile")   (* new: defines scheme() *)

Module("AppLogic")
SchemeStubs("AppLogic")     (* M3 -> Scheme stubs *)

scheme("mathlib")           (* Scheme -> M3 compiled *)
scheme("filters")

ExportSchemeStubs("myapp")
ExportSchemeCompiled("myapp")

importSchemeStubs()
importSchemeCompiled()

implementation("Main")
program("myapp")
\end{lstlisting}

In this example, \texttt{AppLogic.i3} is a Modula-3 interface
whose procedures are exported to Scheme via stubs, while
\texttt{mathlib.scm} and \texttt{filters.scm} are Scheme source
files compiled to native Modula-3 code.  Both kinds of bindings
are registered in the global environment at startup.

%======================================================================
\chapter{The Level~1 Compiler Implementation}
\label{ch:implementation}
%======================================================================

The preceding chapters described the design space and proposed
compilation strategies.  This chapter describes the actual
implementation of the Level~1 compiler: the build system integration,
the test suite, the bugs discovered and fixed during development,
the identifier generation scheme, the measured performance results,
large-scale validation on a production codebase, and a concrete
future optimization that would eliminate virtual dispatch overhead
on the fast path.

The compiler is implemented in \texttt{compiler.scm} (approximately
1000 lines of Scheme), resides in the \texttt{schemecompiler} package
under \texttt{m3-scheme/}, and follows Strategy~B
(\S\ref{sec:stratB}): ahead-of-time compilation from Scheme source to
Modula-3 source, integrated into the CM3 build via a Quake template.

%----------------------------------------------------------------------
\section{Build System Integration}
\label{sec:impl-build}
%----------------------------------------------------------------------

\subsection{The \texttt{scheme()} Quake Template}

The build-time integration is provided by the file
\texttt{schemecompiler/src/scheme.tmpl}, which defines a single
Quake procedure.  The complete template is short enough to reproduce
in full:

\begin{lstlisting}
/* scheme.tmpl -- Quake template for MScheme Level 1 Compiler */

_SCHEME_MSCHEME  = PKG_INSTALL & SL & "mscheme-interactive"
                 & SL & BUILD_DIR & SL & "mscheme"
_SCHEME_COMPILER = PKG_INSTALL & SL & "schemecompiler"
                 & SL & "src" & SL & "compiler.scm"
_SCHEME_EXIT     = "exit"

readonly proc scheme(module_name, scm_file) is
  local src = path_of(scm_file)
  if stale(module_name & ".i3", src)
    local build_scm = module_name & "_scompile.scm"
    > build_scm in
      write("(load \"" & _SCHEME_COMPILER & "\")", CR)
      write("(compile-file \"" & src & "\""
            & " \"" & module_name & "\")", CR)
    end
    _exec([_SCHEME_MSCHEME, build_scm, _SCHEME_EXIT])
  end
  derived_interface(module_name, VISIBLE)
  derived_implementation(module_name)
end

Scheme = scheme
\end{lstlisting}

The template follows the same three-step pattern established by the
\texttt{SchemeStubs} mechanism in \texttt{sstubgen}
(\S\ref{sec:sstubgen}):

\begin{enumerate}
\item \textbf{Staleness check.}  The \texttt{stale()} guard compares
  the modification time of the generated \texttt{.i3} file against
  the Scheme source file (resolved via \texttt{path\_of}).  If the
  output is newer than the input, the compiler is not re-invoked.
  This is the standard CM3 incremental-build idiom.

\item \textbf{Compiler invocation.}  A small temporary Scheme script
  is generated (named \texttt{\emph{module}\_scompile.scm}) that
  loads the compiler and calls \texttt{compile-file} with the source
  path and module name.  The script is executed by the
  \texttt{mscheme-interactive} binary via \texttt{\_exec}.
  This two-stage approach---generate a script, then execute it---avoids
  quoting issues in the Quake-to-Scheme boundary.

\item \textbf{Registration of derived sources.}  The generated
  \texttt{.i3} and \texttt{.m3} files are registered with
  \texttt{derived\_interface} and \texttt{derived\_implementation},
  which tells CM3 to compile them as part of the current package.
  The interface is marked \texttt{VISIBLE}, so other packages can
  import it (e.g., to call \texttt{Install}).
\end{enumerate}

\subsection{Package Structure}

The \texttt{schemecompiler} package's own \texttt{m3makefile}
is minimal:

\begin{lstlisting}
import("cit_util")
template("scheme")
ship_source("compiler.scm")
ship_source("test-input.scm")
interface("SchemeCompiler")
library("schemecompiler")
\end{lstlisting}

The \texttt{template("scheme")} directive makes the
\texttt{scheme.tmpl} file available to downstream packages.  The
\texttt{ship\_source} directives ensure that the compiler source and
the test input file are installed into the package's shipped source
directory (so that the template can locate \texttt{compiler.scm} via
\texttt{PKG\_INSTALL}).  The \texttt{SchemeCompiler.i3} interface is
a placeholder anchor that gives the package a public face in the CM3
namespace.

A client package that wants to compile Scheme source writes:

\begin{lstlisting}
import("schemecompiler")
scheme("CompiledTest", "test-input.scm")
\end{lstlisting}

This single directive compiles all \texttt{define} forms in
\texttt{test-input.scm} to the Modula-3 module \texttt{CompiledTest},
and registers the generated \texttt{CompiledTest.i3} and
\texttt{CompiledTest.m3} for compilation.  The client's
\texttt{Main.m3} then calls:

\begin{lstlisting}[language=Modula3]
CompiledTest.Install(scm);
\end{lstlisting}

\noindent
to register all compiled procedures in the interpreter's global
environment.

\subsection{Differences from the \texttt{SchemeStubs} Pattern}

The \texttt{scheme()} template is simpler than the three-phase
\texttt{SchemeStubs} protocol described in
\S\ref{sec:sstubgen}.  The \texttt{SchemeStubs} mechanism requires
per-interface stub generation (Phase~1), per-package aggregation
(Phase~2), and cross-package aggregation (Phase~3) because stubs must
be registered at module-initialization time, before any user code runs.

The \texttt{scheme()} template requires only a single phase: each
\texttt{scheme()} directive produces one module with an explicit
\texttt{Install} procedure that the client calls at a known point
in its initialization.  This keeps the build integration trivial at
the cost of requiring an explicit \texttt{Install} call.  If
future needs require automatic registration (e.g., for Scheme
libraries used transitively across many packages), the two-phase and
three-phase aggregation patterns from \texttt{SchemeStubs} can be
adopted without changing the compiler itself.

%----------------------------------------------------------------------
\section{Test Suite}
\label{sec:impl-tests}
%----------------------------------------------------------------------

The compiler is validated by a comprehensive correctness test suite
consisting of \textbf{321 test cases} organized into \textbf{12
categories}.  All 321 tests pass on ARM64\_DARWIN (Apple Silicon).

\subsection{Test Structure}

The test suite consists of two files:

\begin{description}
\item[\texttt{test-input.scm}] Contains 80+ Scheme function
  definitions exercising every construct the compiler supports.
  This file is compiled at build time via the \texttt{scheme()}
  directive, producing the \texttt{CompiledTest} module.

\item[\texttt{Main.m3}] A Modula-3 test driver that instantiates an
  MScheme interpreter, calls \texttt{CompiledTest.Install(scm)} to
  register the compiled procedures, and then evaluates 321 test
  expressions via \texttt{scm.loadEvalText()}, comparing each result
  against an expected string representation.
\end{description}

Each test is a call to the procedure \texttt{T}:

\begin{lstlisting}[language=Modula3]
PROCEDURE T(scm: Scheme.T; expr, expected: TEXT)
  RAISES {Scheme.E} =
VAR got := SchemeUtils.StringifyT(scm.loadEvalText(expr));
BEGIN
  INC(total);
  IF Text.Equal(got, expected) THEN
    INC(pass)
  ELSE
    INC(fail);
    (* report failure *)
  END
END T;
\end{lstlisting}

The key point is that the compiled procedures are called through
the normal interpreter evaluation path: \texttt{loadEvalText}
parses the expression, the evaluator resolves the function name
via the global environment, and dispatches through the standard
\texttt{SchemeProcedure.T.apply} interface.  The compiled procedures
are indistinguishable from interpreted ones at the call site.  This is
a consequence of the design decision to implement compiled procedures
as \texttt{SchemeProcedure.T} subtypes (\S\ref{sec:proctype}).

\subsection{Test Categories}

\begin{center}
\begin{longtable}{lcp{7.5cm}}
\toprule
\textbf{Category} & \textbf{Tests} & \textbf{Constructs exercised} \\
\midrule
\endhead
Arithmetic & 54 &
  \texttt{factorial}, \texttt{fibonacci}, \texttt{square}, \texttt{cube},
  \texttt{abs-val}, \texttt{max-of-two}, \texttt{min-of-two},
  \texttt{power}, \texttt{sum-range}, \texttt{even?}, \texttt{odd?},
  \texttt{gcd}, \texttt{negate}, \texttt{double}, \texttt{triple},
  \texttt{add3}, \texttt{add4}.
  Exercises \texttt{if}, recursion, multi-arity calls (1--4 args),
  free-variable access to primitives (\texttt{+}, \texttt{*},
  \texttt{-}, \texttt{=}, \texttt{<}, \texttt{>}, \texttt{quotient},
  \texttt{remainder}). \\[0.5em]
List Operations & 56 &
  \texttt{my-length}, \texttt{sum-list}, \texttt{product-list},
  \texttt{my-last}, \texttt{my-nth}, \texttt{my-append},
  \texttt{my-reverse-acc}, \texttt{my-map}, \texttt{my-filter},
  \texttt{my-member}, \texttt{count-if}, \texttt{take}, \texttt{drop},
  \texttt{zip}, \texttt{flatten}, \texttt{list-ref-safe},
  \texttt{iota-acc}, \texttt{repeat-val}.
  Exercises \texttt{null?}, \texttt{car}, \texttt{cdr}, \texttt{cons},
  \texttt{pair?}, quoted empty lists, higher-order calls passing
  compiled functions as arguments. \\[0.5em]
Cond & 17 &
  \texttt{sign}, \texttt{classify}, \texttt{fizzbuzz},
  \texttt{letter-grade}.
  Exercises multi-clause \texttt{cond} with \texttt{else},
  numeric and string return values, nested comparisons. \\[0.5em]
Begin & 4 &
  \texttt{begin-test}, \texttt{begin-with-side-effect}.
  Exercises sequencing with discarded intermediate results,
  temporary variable generation. \\[0.5em]
Let / Let* & 18 &
  \texttt{hypotenuse}, \texttt{circle-area}, \texttt{let-nested},
  \texttt{let-star-chain}, \texttt{let-star-deps},
  \texttt{swap-via-let}, \texttt{multi-let}.
  Exercises parallel binding (\texttt{let}), sequential binding
  (\texttt{let*}), multi-variable \texttt{let}, nested \texttt{let},
  initializer dependencies in \texttt{let*}. \\[0.5em]
And / Or / Not & 30 &
  \texttt{both-positive?}, \texttt{either-zero?}, \texttt{in-range?},
  \texttt{out-of-range?}, \texttt{three-and}, \texttt{three-or},
  \texttt{not-zero?}, \texttt{xor-bool}.
  Exercises short-circuit evaluation, boolean return values,
  multi-operand \texttt{and}/\texttt{or}, nested boolean logic. \\[0.5em]
Set! (Mutation) & 9 &
  \texttt{set-param}, \texttt{set-in-let}, \texttt{set-accumulate}.
  Exercises mutation of parameters and let-bound variables. \\[0.5em]
Self-Tail-Call & 22 &
  \texttt{loop-sum}, \texttt{count-down}, \texttt{tail-factorial},
  \texttt{tail-fib-iter}, \texttt{tail-length},
  \texttt{tail-sum-list}, \texttt{tail-last}, \texttt{tail-min},
  \texttt{tail-max}, \texttt{repeat-string}.
  Exercises the LOOP optimization for functions with 1--3
  parameters, including \texttt{if} and \texttt{cond} as the
  branch construct, and string accumulation. \\[0.5em]
Constants & 14 &
  \texttt{return-true}, \texttt{return-false}, \texttt{return-zero},
  \texttt{return-one}, \texttt{return-string},
  \texttt{return-empty-list}, \texttt{return-symbol},
  \texttt{return-quoted-list}, \texttt{return-char},
  \texttt{identity}, \texttt{constant-42},
  \texttt{return-pi-approx}, \texttt{return-negative}.
  Exercises every constant type: booleans, integers,
  strings, the empty list, symbols, quoted lists, characters.
  Includes zero-argument functions. \\[0.5em]
Recursive Algorithms & 38 &
  \texttt{ackermann}, \texttt{hanoi-moves}, \texttt{depth},
  \texttt{tree-sum}, \texttt{tree-count}, \texttt{collatz-steps},
  \texttt{digit-sum}, \texttt{num-digits},
  \texttt{is-palindrome-helper}, \texttt{list-equal?}.
  Exercises deep non-tail recursion, tree traversal, mutual
  recursion between compiled functions, \texttt{cond} with
  compound boolean tests (\texttt{and}/\texttt{or} in test
  position). \\[0.5em]
Higher-Order & 4 &
  \texttt{apply-twice}, \texttt{compose-apply}.
  Exercises passing compiled functions as arguments (free variables
  of procedure type), calling function-valued parameters. \\[0.5em]
Mixed Constructs & 16 &
  \texttt{complex-cond-let}, \texttt{nested-if-let},
  \texttt{and-or-mix}, \texttt{multi-begin-let},
  \texttt{cond-with-begin}, \texttt{if-chain}.
  Exercises deep nesting of \texttt{cond}/\texttt{let}/\texttt{let*}/\texttt{begin}/\texttt{if},
  boolean expressions in non-tail position. \\[0.5em]
Edge Cases & 19 &
  \texttt{zero-args}, \texttt{single-arg}, \texttt{five-args},
  \texttt{deeply-nested}, \texttt{many-lets}, \texttt{let-shadow},
  \texttt{cond-many}, \texttt{boolean-identity}.
  Exercises 0, 1, and 5 parameter functions; 5-level nested
  \texttt{if}; 4-level nested \texttt{let}; variable shadowing
  in \texttt{let}; \texttt{cond} with 5 clauses; returning literal
  \texttt{\#t} and \texttt{\#f}. \\[0.5em]
\bottomrule
\end{longtable}
\end{center}

\subsection{Benchmark Harness}

Following the 321 correctness tests, the test program includes a
performance benchmark section.  Four benchmarks are timed in two
configurations: first with the compiled procedures (as installed by
\texttt{CompiledTest.Install}), then again after redefining the same
functions as interpreted closures via \texttt{loadEvalText}.

The benchmark harness uses \texttt{Time.Now()} for wall-clock timing.
Each benchmark is executed once (no warmup iterations), which is
acceptable because the workloads are large enough to dominate any
startup overhead.  The \texttt{loadEvalText} call includes parsing
overhead; for the self-tail-call benchmark (\texttt{loop-sum}), this
parsing time dominates the measurement, as discussed in
\S\ref{sec:impl-perf}.

%----------------------------------------------------------------------
\section{Bugs Found and Fixed}
\label{sec:impl-bugs}
%----------------------------------------------------------------------

Development of the test suite and large-scale validation
(\S\ref{sec:impl-cspc}) exposed five compiler bugs.  The first two
were subtle interactions between compilation assumptions and Modula-3
scoping semantics; the remaining three were discovered during
stress testing on the CSP compiler.

\subsection{The \texttt{if} Sentinel Bug}
\label{sec:bug-if}

\paragraph{Symptom.}
The test \texttt{(boolean-identity \#f)} entered an infinite loop
instead of returning \texttt{\#f}.

\paragraph{Root cause.}
The compiler used the Scheme value \texttt{\#f} as both the boolean
false literal \emph{and} the sentinel for ``no else branch present.''
When the \texttt{if} special form was parsed, the compiler checked
\texttt{(pair?\ (cdddr form))} to determine whether an else branch
existed; if not, it passed \texttt{\#f} as the else expression.
Later, \texttt{gen-if} tested \texttt{(not alt)} to decide whether
to emit an \texttt{ELSE} clause.

When the source code was literally \texttt{(if x \#t \#f)}, the
parsed else expression was the boolean \texttt{\#f}.  The
\texttt{(not \#f)} test returned true, so the compiler concluded
there was no else branch and omitted the \texttt{ELSE} clause.
The generated code for a self-tail-call function became:

\begin{lstlisting}[language=Modula3]
LOOP
  IF SchemeBoolean.TruthO(a_x) THEN
    RETURN SchemeBoolean.True()
  END;
  (* no ELSE, no RETURN -- falls through to next LOOP iteration *)
END
\end{lstlisting}

When the argument was false, the \texttt{IF} body was not entered,
and the \texttt{LOOP} iterated forever.

\paragraph{Fix.}
Replace the sentinel with a unique symbol:

\begin{lstlisting}[language=Scheme]
;; In gen-expr, the if case:
((eq? (car expr) 'if)
 (gen-if (cadr expr) (caddr expr)
         (if (pair? (cdddr expr))
             (cadddr expr)
             '*no-alt*)     ; was: #f
         target ctx depth))

;; In gen-if:
(if (not (eq? alt '*no-alt*))  ; was: (not alt)
    ...)
\end{lstlisting}

The symbol \texttt{'*no-alt*} is never a valid Scheme expression (it
would evaluate to an unbound variable), so there is no ambiguity.
The sentinel is tested with \texttt{eq?}, a pointer comparison, so
the performance impact is nil.

\subsection{The \texttt{let}/\texttt{let*} Shadowing Bug}
\label{sec:bug-let}

\paragraph{Symptom.}
The test \texttt{(let-shadow 5)} returned \texttt{0} instead of
\texttt{30}.  The function:

\begin{lstlisting}[language=Scheme]
(define (let-shadow x)
  (let ((x (+ x 10)))
    (let ((x (* x 2)))
      x)))
\end{lstlisting}

With $x = 5$, should compute: inner $x = (5 + 10) = 15$, outer
$x = 15 \times 2 = 30$.

\paragraph{Root cause.}
The original \texttt{gen-let} implementation declared the let-bound
variable and evaluated its initializer in the \emph{same} Modula-3
scope:

\begin{lstlisting}[language=Modula3]
(* BUGGY generated code for (let ((x (+ x 10))) ...) *)
VAR a_x : SchemeObject.T; BEGIN
  a_x := (* evaluate (+ x 10) *)
    ...
  (* but the reference to x in the initializer
     resolves to a_x, which is uninitialized! *)
END;
\end{lstlisting}

In Modula-3, the \texttt{VAR} declaration introduces \texttt{a\_x}
into scope for the entire block, including the initializer
expression (Nelson~[5, Chapter~2]).  When the let-bound variable
shadows a parameter with the same name, the initializer reads the
uninitialized inner variable instead of the outer parameter.  The
variable is initialized to zero by default (or \texttt{NIL} for
traced references), hence the incorrect result.

Scheme's \texttt{let} semantics require that all initializer
expressions be evaluated in the \emph{enclosing} scope, before any
of the new bindings are visible.  This is explicitly stated in R5RS:
``The \texttt{<init>}s are evaluated in the current environment (in
some unspecified order), the \texttt{<variable>}s are bound to fresh
locations holding the results, and the \texttt{<body>} is evaluated
in the extended environment.''

\paragraph{Fix.}
The fix uses a two-phase approach: evaluate initializers into
temporaries in the outer scope, then declare the let-bound variables
in a nested scope initialized from the temporaries:

\begin{lstlisting}[language=Scheme]
(define (gen-let bindings body target ctx depth)
  (let* ((bvars (map car bindings))
         (temps (map (lambda (b) (fresh-temp ctx))
                     bindings))
         (new-ctx ...))
    (string-append
     ;; Phase 1: temps in outer scope
     (indent depth) (L "VAR")
     ... ;; declare temps
     (indent depth) (L "BEGIN")
     ... ;; evaluate inits into temps (using ctx,
         ;; not new-ctx, so x refers to the param)

     ;; Phase 2: let-bound vars in nested scope
     (indent (+ depth 1)) (L "VAR")
     ... ;; declare a_x := temp
     (indent (+ depth 1)) (L "BEGIN")
     ... ;; body using new-ctx
     (indent (+ depth 1)) (L "END;")
     (indent depth) (L "END;"))))
\end{lstlisting}

The generated Modula-3 now has the correct scoping:

\begin{lstlisting}[language=Modula3]
(* CORRECT generated code for (let ((x (+ x 10))) ...) *)
VAR
  t_0 : SchemeObject.T;
BEGIN
  (* evaluate initializer in outer scope:
     a_x here refers to the parameter *)
  t_0 := SchemeLongReal.Add(a_x, self.k_0);
  VAR
    a_x : SchemeObject.T := t_0;
  BEGIN
    (* body: a_x here refers to the let-bound var *)
    ...
  END;
END;
\end{lstlisting}

The same fix was applied to \texttt{gen-let*}, where the sequential
scoping of \texttt{let*} requires each binding to be evaluated before
the next variable is introduced:

\begin{lstlisting}[language=Modula3]
(* (let* ((a (+ x 1)) (b (* a 2))) ...) *)
VAR t_0 : SchemeObject.T; BEGIN
  t_0 := SchemeLongReal.Add(a_x, self.k_0);
  VAR a_a : SchemeObject.T := t_0; BEGIN
    VAR t_1 : SchemeObject.T; BEGIN
      t_1 := SchemeLongReal.Mul(a_a, self.k_1);
      VAR a_b : SchemeObject.T := t_1; BEGIN
        ...
      END;
    END;
  END;
END;
\end{lstlisting}

\paragraph{Temp counter sharing.}
An additional subtlety: the original code created a fresh
compilation context via \texttt{make-ctx} for the inner scope,
which initialized a new mutable counter cell \texttt{(list 0)} for
temporary variable names.  This caused temporaries in the inner
scope to collide with those in the outer scope (both starting at
\texttt{t\_0}).  The fix shares the parent context's counter cell:

\begin{lstlisting}[language=Scheme]
(define (gen-let bindings body target ctx depth)
  (let* (...
         (new-ctx (list (ctx-name ctx)
                        (append bvars (ctx-params ctx))
                        (ctx-free-vars ctx)
                        (ctx-constants ctx)
                        (car (cddddr ctx))  ; SHARE counter
                        (ctx-self-tail? ctx))))
    ...))
\end{lstlisting}

The expression \texttt{(car (cddddr ctx))} extracts the mutable
counter cell (a one-element list holding the current count) from the
parent context and installs it in the child context, so that
\texttt{fresh-temp} increments the same counter in both scopes.

\subsection{The Bodyless \texttt{define} Crash}
\label{sec:bug-bodyless}

\paragraph{Symptom.}
The compiler crashed with an error on \texttt{(define (bigint-dbg .\ x))}
--- a definition with no body, used in the CSP compiler as a debugging
stub.

\paragraph{Root cause.}
The compilability checker called \texttt{(cdddr form)} to access the
body of a define form, which fails when the form has only two elements
(name and parameter list, no body).  Similarly, the top-level filter
that extracts compilable defines assumed every \texttt{define} form
had a body.

\paragraph{Fix.}
Guard both \texttt{define-compilable?}\ and the \texttt{all-defines}
filter with \texttt{(pair?\ (cddr form))} and
\texttt{(pair?\ (cddr f))} respectively, so that bodyless defines are
silently skipped rather than crashing.

\subsection{The Short \texttt{if} Crash}
\label{sec:bug-short-if}

\paragraph{Symptom.}
The compiler crashed on \texttt{(if test)} --- an \texttt{if} form
with only a test expression and no consequent or alternative.

\paragraph{Root cause.}
The \texttt{expr-compilable?}\ function called \texttt{(cdddr expr)}
unconditionally on \texttt{if} forms to check for an alternative
branch.  When the form had fewer than four elements, this raised an
error.

\paragraph{Fix.}
Guard the alternative-branch check with \texttt{(pair?\ (cddr expr))},
so that degenerate \texttt{if} forms are rejected by the compilability
checker rather than crashing.

\subsection{The \texttt{collect-constants} Scanning Bug}
\label{sec:bug-constants}

\paragraph{Symptom.}
After compiling the CSP compiler's 587 definitions, the Modula-3
build produced 12 errors of the form ``\texttt{unknown qualification
`.'\,(k\_0)}''.  The generated code referenced constant fields
(\texttt{self.k\_0}) that did not exist in the object type declaration.

\paragraph{Root cause.}
The \texttt{collect-constants} function, which walks the AST to find
all quoted symbols, strings, and other constants that need to be
stored as object fields, had a bug in its recursive traversal.
The else branch (for general list expressions) used
\texttt{(cdr expr)} instead of \texttt{expr}:

\begin{lstlisting}[language=Scheme]
;; BUGGY: skips first element of every sub-list
(else (apply append (map collect-constants (cdr expr))))
\end{lstlisting}

This skipped the first element of every list, which meant constants
appearing as the first element of a sub-expression---such as the
quoted symbol \texttt{'+inf} in \texttt{(member?\ '+inf blist)}---were
never collected.  The code generator emitted references to these
constants as \texttt{self.k\_$n$}, but the type declaration had no
corresponding field because the constant was never registered.

\paragraph{Fix.}
Change \texttt{(cdr expr)} to \texttt{expr} so that all elements
of every sub-expression are scanned:

\begin{lstlisting}[language=Scheme]
;; FIXED: scans all elements including the first
(else (apply append (map collect-constants expr)))
\end{lstlisting}

As an additional safety net, the \texttt{constant-ref} function was
extended with fallback inline generation: if a constant is not found
in the constants list (which should never happen after the fix), it
generates the constant value inline rather than emitting a dangling
field reference.  This makes the compiler robust against future
scanning omissions.

%----------------------------------------------------------------------
\section{Identifier Generation}
\label{sec:impl-idents}
%----------------------------------------------------------------------

A critical design decision in the compiler is how Scheme identifiers
(which may contain characters like \texttt{?}, \texttt{!}, \texttt{->},
and \texttt{+}) are mapped to valid Modula-3 identifiers.

\subsection{The Name Collision Problem}

An early version of the compiler used a transliteration function
\texttt{m3-ident} that mapped each Scheme character to a valid
Modula-3 identifier character (e.g., \texttt{?}~$\to$~\texttt{Q},
\texttt{!}~$\to$~\texttt{B}, \texttt{-}~$\to$~\texttt{\_}).
Scheme names were prefixed by their role (\texttt{a\_} for
parameters, \texttt{b\_} for bindings) and then transliterated:
\texttt{null?} became \texttt{a\_nullQ}, \texttt{set-car!} became
\texttt{b\_set\_carB}.

This scheme had a fundamental problem: user-provided Scheme
identifiers directly influenced the generated Modula-3 identifiers,
creating the possibility of name collisions.  Two different Scheme
names could map to the same Modula-3 identifier if their
transliterations coincided.  More importantly, a let-bound variable
could shadow a parameter with the same transliterated name in
unexpected ways when the Scheme names differed but the Modula-3
names matched.  For a compiler generating code for arbitrary user
programs, this was unacceptable.

\subsection{Serial-Numbered Identifiers}

The solution is to generate all Modula-3 identifiers using serial
numbers, with no dependence on the original Scheme names.  The
naming conventions are:

\begin{center}
\begin{tabular}{lll}
\toprule
\textbf{Prefix} & \textbf{Role} & \textbf{Example} \\
\midrule
\texttt{p}     & Procedure types/names & \texttt{p0}, \texttt{p1}, \texttt{p2} \\
\texttt{a\_}   & Parameters            & \texttt{a\_0}, \texttt{a\_1}, \texttt{a\_2} \\
\texttt{b\_}   & Free-variable bindings & \texttt{b\_0}, \texttt{b\_1}, \texttt{b\_2} \\
\texttt{v\_}   & Let/let*-bound variables & \texttt{v\_0}, \texttt{v\_1}, \texttt{v\_2} \\
\texttt{k\_}   & Constants              & \texttt{k\_0}, \texttt{k\_1}, \texttt{k\_2} \\
\texttt{t\_}   & Temporaries            & \texttt{t\_0}, \texttt{t\_1}, \texttt{t\_2} \\
\bottomrule
\end{tabular}
\end{center}

Each compiled define receives a global serial number: the first
define in a file becomes \texttt{p0}, the second \texttt{p1}, and
so on.  Within each define, parameters and free variables are
assigned indices by their declaration order.  Let-bound variables
use a per-function counter that is shared across nested scopes (as
described in \S\ref{sec:bug-let}) to prevent collisions.

The compiler maintains the mapping from Scheme names to serial-numbered
Modula-3 names in association lists (\texttt{param-map} and
\texttt{fv-map}) stored in the compilation context.  When
\texttt{var-ref} needs to generate code for a variable reference,
it looks up the Scheme symbol in these maps:

\begin{lstlisting}[language=Scheme]
(define (var-ref sym ctx)
  (let ((p (assq sym (ctx-param-map ctx))))
    (if p (cdr p)
        (let ((f (assq sym (ctx-fv-map ctx))))
          (if f (string-append "self." (cdr f) ".get()")
              (error "unresolved variable" sym))))))
\end{lstlisting}

This guarantees that no user-provided string ever appears in a
generated Modula-3 identifier.  The only place Scheme names appear
in the generated code is in string literals passed to
\texttt{SchemeSymbol.Symbol()} during registration, where they are
harmless.

\subsection{Self-Tail-Call and Let Shadowing}

The serial-numbered scheme interacts cleanly with the self-tail-call
optimization.  In the generated \texttt{LOOP}, the parameter variables
are always the original \texttt{a\_0}, \texttt{a\_1}, etc.  When a
let-bound variable shadows a parameter, it gets a distinct name from
the \texttt{v\_} series, so the tail-call reassignment
\texttt{a\_0~:=~\ldots;\ a\_1~:=~\ldots} always refers to the
correct LOOP variables, even from within a nested let scope that has
shadowed one of the parameters.

\subsection{Define Deduplication}
\label{sec:impl-dedup}

When compiling a large Scheme codebase, it is common for the same
file to be loaded more than once (e.g., a utility library loaded
by multiple modules that are later concatenated).  This produces
duplicate \texttt{define} forms that would cause ``symbol
already defined'' errors in the generated Modula-3.

The compiler applies a deduplication pass before compilation:
\texttt{deduplicate-defines} walks the list of extracted defines
in reverse order, keeping only the \emph{first} occurrence of each
function name (which corresponds to the \emph{last} definition in
the source, preserving Scheme's redefinition semantics).  Duplicate
definitions are silently dropped with a diagnostic message.

In the CSP compiler test (\S\ref{sec:impl-cspc}), this removed 4
duplicate definitions from the 867 extracted defines, reducing the
unique count to 863.

%----------------------------------------------------------------------
\section{Performance Results}
\label{sec:impl-perf}
%----------------------------------------------------------------------

The test program includes a benchmark section that compares compiled
and interpreted execution of four workloads.  The benchmarks were
run on ARM64\_DARWIN (Apple M1, macOS).

\begin{center}
\begin{tabular}{lccc}
\toprule
\textbf{Benchmark} & \textbf{Compiled} & \textbf{Interpreted} & \textbf{Speedup} \\
\midrule
\texttt{(fibonacci 25)}       & 225 ms    & 2,324 ms   & 10.3$\times$ \\
\texttt{(fibonacci 30)}       & 5,603 ms  & 24,092 ms  & 4.3$\times$ \\
\texttt{(loop-sum 10000 0)}   & 36 ms     & 83 ms      & 2.3$\times$ \\
\texttt{(ackermann 3 4)}      & 13 ms     & 81 ms      & 6.4$\times$ \\
\bottomrule
\end{tabular}
\end{center}

\subsection{Analysis}

The compiled code delivers a \textbf{4--10$\times$ speedup} on
recursive workloads.  The speedup varies by benchmark for several
reasons:

\paragraph{Fibonacci.}
The \texttt{fibonacci} benchmark is the purest test of call
overhead.  Each call performs a comparison, two subtractions, two
recursive calls, and one addition---very little per-call work.
The interpreter's per-call overhead (type dispatch,
special-form recognition, operator TYPECASE, and argument-list
construction, as analyzed in \S\ref{ch:analysis}) dominates.  The
compiled code eliminates all of this except the \texttt{NARROW} and
virtual dispatch through \texttt{SchemeProcedure.T.apply1}.  The
10.3$\times$ speedup on \texttt{(fibonacci 25)} reflects this.  The
lower 4.3$\times$ on \texttt{(fibonacci 30)} is a measurement
artifact: the longer absolute runtime amortizes fixed overhead.

\paragraph{Ackermann.}
The Ackermann function involves
deeply nested non-tail recursion with very small bodies.  The
6.4$\times$ speedup is consistent with the fibonacci result: per-call
overhead dominates, and the compiler eliminates most of it.

\paragraph{Loop-sum.}
The \texttt{loop-sum} benchmark exercises the self-tail-call
optimization: the function is compiled to a Modula-3 \texttt{LOOP}
with parameter reassignment, eliminating all call overhead.  The
\texttt{LOOP} body executes a single comparison, a subtraction,
and an addition---on ARM64 with the CM3 C backend, this compiles
to a handful of machine instructions per iteration.  For
$n = 10{,}000$, the compiled \texttt{LOOP} body executes in
microseconds.

The measured 2.3$\times$ speedup appears surprisingly modest for
a benchmark where call overhead is eliminated entirely.  The
explanation is the \texttt{loadEvalText} overhead in the benchmark
harness.  Each benchmark invocation:

\begin{enumerate}
\item Parses the expression string \texttt{"(loop-sum 10000 0)"},
  constructing an AST.
\item Evaluates the AST: looks up \texttt{loop-sum} in the global
  environment, evaluates the arguments \texttt{10000} and \texttt{0},
  constructs an argument list, and dispatches through
  \texttt{apply2}.
\item The compiled \texttt{LOOP} executes.
\item The result is returned.
\end{enumerate}

Steps 1--2 are identical for both compiled and interpreted
configurations.  For the compiled configuration, step~3 is
near-instantaneous.  For the interpreted configuration, step~3
involves 10,000 recursive calls through the eval loop, each with
the full per-call overhead.  The 36~ms measured for ``compiled''
is almost entirely step~1 (parsing) and step~2 (eval-loop entry),
not step~3 (the actual computation).

\subsection{The Bottleneck: Binding Lookup and Virtual Dispatch}

For non-self-tail-call functions, the remaining per-call overhead
in compiled code is:

\begin{enumerate}
\item \texttt{self.b\_\emph{foo}.get()} --- one binding lookup
  (pointer dereference, array index, load) to fetch the callee.
\item \texttt{NARROW(\emph{proc}, SchemeProcedure.T)} --- a runtime
  type check.
\item \texttt{.apply1(interp, arg)} --- a virtual method dispatch
  (indirect call through the method table).
\end{enumerate}

On ARM64, this sequence compiles to approximately 5--8 instructions
(two loads for the binding, one branch for the type check, one
indirect branch for the virtual call).  This is far cheaper than the
interpreter's per-call overhead but still measurable for
call-intensive benchmarks like \texttt{fibonacci}.

The self-tail-call optimization eliminates this overhead entirely
for tail-recursive functions by replacing the call with a
\texttt{LOOP}/parameter-reassignment pattern that compiles to a
backward branch.

%----------------------------------------------------------------------
\section{Future Optimization: Intra-Module Direct Calls}
\label{sec:impl-future}
%----------------------------------------------------------------------

The Level~1 design (\S\ref{sec:redefinition}) mandates that all calls
to free variables go through the binding cell, preserving redefinition
semantics.  However, for calls between functions \emph{in the same
compilation unit}, a hybrid approach can eliminate virtual dispatch
overhead on the fast path while preserving full redefinition semantics
on the slow path.

\subsection{The Direct-Call Optimization}

For each function \texttt{foo} defined in the compilation unit, the
\texttt{Install} procedure already holds a direct typed reference
\texttt{obj\_foo : Compiled\_foo}.  This reference can be stored
as an additional field on each compiled procedure object that calls
\texttt{foo}:

\begin{lstlisting}[language=Modula3]
TYPE Compiled_bar = SchemeProcedure.T OBJECT
    b_foo : SchemeEnvironmentBinding.T;  (* binding *)
    d_foo : Compiled_foo;  (* direct ref, or NIL *)
    ...
  END;
\end{lstlisting}

At each call site for \texttt{foo}, the generated code performs a
pointer comparison:

\begin{lstlisting}[language=Modula3]
VAR proc := self.b_foo.get();
BEGIN
  IF proc = self.d_foo THEN
    (* fast path: direct call, no NARROW, no vtable *)
    RETURN Apply1_foo(self.d_foo, interp, arg)
  ELSE
    (* slow path: general dispatch *)
    RETURN NARROW(proc, SchemeProcedure.T).apply1(interp, arg)
  END
END
\end{lstlisting}

\subsection{Why This Preserves Redefinition Semantics}

The optimization always checks the live binding first.  The pointer
comparison \texttt{proc = self.d\_foo} asks: ``is the value currently
stored in \texttt{foo}'s binding cell the same object that was
installed at registration time?''

\begin{itemize}
\item \textbf{If yes:} The callee has not been redefined.  The
  direct call is safe and semantically equivalent to the general
  dispatch---but faster, because it avoids \texttt{NARROW} and
  the virtual method dispatch.

\item \textbf{If no:} The callee has been redefined (either from
  the REPL, via \texttt{set!}, or via a re-\texttt{define}).
  The code falls back to the general dispatch path, which correctly
  calls whatever is now in the binding cell---compiled, interpreted,
  or a completely different function.
\end{itemize}

The pointer comparison is a single ARM64 \texttt{cmp} instruction
followed by a conditional branch.  Modern branch predictors will
predict the fast path with near-100\% accuracy because redefinition
is rare in normal execution.  The fast path then avoids both the
\texttt{NARROW} check (which involves a type tag comparison) and
the virtual dispatch (which involves an indirect load from the
method table and an indirect branch).

\subsection{Scope and Limitations}

This optimization applies only to calls between functions defined in
the same \texttt{scheme()} compilation unit.  For calls to functions
defined elsewhere (primitives, functions from other packages), the
compiler has no direct typed reference to store, and the general
binding-lookup path is used unconditionally.

The optimization also does not help with calls through
function-valued parameters (higher-order calls), where the callee
is never known at compile time.

For a benchmark like \texttt{fibonacci}---where both the recursive
calls are to the same function in the same compilation unit---this
optimization would eliminate the remaining per-call overhead and
approach the performance of hand-written Modula-3 code.  The
optimization has not yet been implemented.

%----------------------------------------------------------------------
\section{Large-Scale Validation: The CSP Compiler}
\label{sec:impl-cspc}
%----------------------------------------------------------------------

The unit test suite (\S\ref{sec:impl-tests}) validates correctness on
isolated functions.  To validate the compiler on a real, large-scale
Scheme codebase, we compiled the entire CSP (Communicating Sequential
Processes) compiler from the Intel Labs async-toolkit project.  The
CSP compiler (\texttt{cspc}) is a production tool for asynchronous
circuit design, written in MScheme, totaling approximately 17,000
lines of Scheme across $\sim$50 source files.

\subsection{Compilation Results}

The 50 source files were concatenated into a single input file
(16,313 lines).  The compiler extracted 867 top-level
\texttt{define} forms.  After deduplication (\S\ref{sec:impl-dedup}),
863 unique definitions remained.  Of these:

\begin{center}
\begin{tabular}{lr}
\toprule
\textbf{Category} & \textbf{Count} \\
\midrule
Successfully compiled    & 587 (68\%) \\
Skipped (unsupported constructs) & 276 (32\%) \\
\midrule
Total unique defines     & 863 \\
Duplicates removed       & 4 \\
\bottomrule
\end{tabular}
\end{center}

The 276 skipped definitions break down by blocking construct:

\begin{center}
\begin{tabular}{lc}
\toprule
\textbf{Blocking construct} & \textbf{Count} \\
\midrule
\texttt{lambda} (in \texttt{map}/\texttt{filter}/\texttt{for-each} args) & $\sim$200 \\
Named \texttt{let}         & $\sim$50 \\
Rest parameters (\texttt{.\ args}) & $\sim$15 \\
\texttt{do} loops          & $\sim$5 \\
Other (internal define, letrec, etc.) & $\sim$6 \\
\bottomrule
\end{tabular}
\end{center}

As with the unit tests, \texttt{lambda} is by far the dominant
gap.  Adding lambda support alone would push the compilation ratio
well above 80\%.

\subsection{Correctness Verification}

To verify that the compiled definitions produce identical results
to the interpreted originals, the compiled CSP compiler was run on
four MiniMIPS processor designs: Fetch, Decode, Execute, and
RegFile.  These are real circuit specifications that exercise
the full depth of the CSP compiler's analysis, transformation, and
code generation passes.

The output of the compiled configuration (587 compiled
+ 276 interpreted definitions) was compared against the output of
the fully interpreted configuration (all 863 definitions interpreted).
The results were \textbf{byte-for-byte identical} on all four
test cases.

This is strong evidence that the compiler preserves semantics
faithfully: the 587 compiled definitions interoperate correctly
with each other, with the 276 interpreted definitions, and with
the MScheme runtime primitives.

\subsection{Bugs Discovered}

The CSP compiler stress test uncovered three additional bugs
(\S\ref{sec:bug-bodyless}, \S\ref{sec:bug-short-if},
\S\ref{sec:bug-constants}) that the unit test suite did not expose,
because the unit tests used only well-formed, conventional Scheme
code.  The CSP compiler, being a large production codebase,
contained edge cases---bodyless debugging stubs, degenerate
\texttt{if} forms, and constants in unusual AST positions---that
exercised code paths the unit tests missed.

This validates the importance of testing on real codebases in
addition to synthetic test suites.

\subsection{Current Restrictions}

The Level~1 compiler supports the following Scheme constructs in
compiled definitions:

\begin{center}
\begin{tabular}{ll}
\toprule
\textbf{Supported} & \textbf{Not yet supported} \\
\midrule
\texttt{if}, \texttt{cond} (with \texttt{else})
  & \texttt{lambda} \\
\texttt{let}, \texttt{let*}
  & Named \texttt{let} \\
\texttt{begin}
  & \texttt{letrec} \\
\texttt{and}, \texttt{or}, \texttt{not}
  & Internal \texttt{define} \\
\texttt{set!}
  & \texttt{cond} with \texttt{=>} \\
\texttt{quote} (all types)
  & \texttt{case}, \texttt{do} \\
Self-tail-call optimization
  & Rest parameters \\
1--4 fixed parameters
  & General tail-call optimization \\
Free-variable capture via bindings
  & \texttt{quasiquote} \\
Constants (numbers, strings, symbols,
  & \texttt{call/cc} \\
\quad booleans, chars, null, quoted lists)
  & \\
\bottomrule
\end{tabular}
\end{center}

Definitions using any unsupported construct are silently skipped
by the compilability checker and fall back to the interpreter with
no loss of semantics.  The compiled and interpreted definitions
interoperate transparently through the
\texttt{SchemeProcedure.T} interface (\S\ref{sec:proctype}).

%======================================================================
\chapter{Remaining Compiler Gaps and How to Close Them}
\label{ch:gaps}
%======================================================================

The compiler currently supports a productive subset of Scheme: all
of the constructs used in ordinary first-order programming are
compiled to efficient Modula-3 code.  Definitions that use
unsupported constructs are silently skipped by the compilability
checker (\S\ref{ch:boundary}), falling back to the
interpreter with no loss of semantics.

This chapter catalogues the remaining gaps---both in language
coverage and in call optimization---and describes concrete
strategies for closing each one.  The gaps are ordered roughly
by impact, measured by how many definitions in the compiler's own
source code (\texttt{compiler.scm}) each gap prevents from being
compiled.  Self-compilation is a useful stress test.  In
addition, the large-scale CSP compiler validation
(\S\ref{sec:impl-cspc}) provides empirical data on which
constructs matter most in a production codebase.

%----------------------------------------------------------------------
\section{Missing Language Constructs}
\label{sec:gaps-lang}
%----------------------------------------------------------------------

\subsection{\texttt{lambda} --- Closures}
\label{sec:gap-lambda}

\paragraph{Impact.}
This is the single largest gap.  In \texttt{compiler.scm},
31 definitions contain \texttt{lambda} expressions---nearly all as
arguments to \texttt{map} and \texttt{filter}---and all 31 are
skipped.  Closing this gap alone would bring the self-compilation
ratio from 38/59 to 50/59.

\paragraph{What is needed.}
A \texttt{lambda} expression creates a first-class closure: an
anonymous procedure that captures the lexical environment at its
point of definition.  The compiler must:

\begin{enumerate}
\item Generate a \texttt{SchemeProcedure.T} subclass for the lambda
  body, just as it does for a top-level \texttt{define}.

\item At the point where the lambda appears, allocate a \texttt{NEW}
  instance of that subclass and populate its free-variable bindings.

\item The lambda's free variables include both the enclosing
  function's parameters (available as local Modula-3 variables) and
  free variables inherited from the enclosing scope (available as
  binding cells on \texttt{self}).
\end{enumerate}

\paragraph{Strategy.}
The compilation of a \texttt{lambda} is structurally identical to
the compilation of a top-level \texttt{define}: analyze the body
for free variables, generate a procedure object type, and emit
\texttt{apply} methods.  The difference is that the object is
allocated inline (within the enclosing function's generated code)
rather than in the \texttt{Install} procedure.

For free variables that are parameters of the enclosing function,
the lambda closure must capture them by \emph{value} (not by
binding cell), since parameters are Modula-3 local variables.
This can be done by adding \texttt{SchemeObject.T} fields to the
lambda's object type, initialized at allocation time:

\begin{lstlisting}[language=Modula3]
(* (lambda (y) (+ x y)) where x is a parameter *)
TYPE Lambda_0 = SchemeProcedure.T OBJECT
    captured_x : SchemeObject.T;  (* captured from enclosing scope *)
    b_plus     : SchemeEnvironmentBinding.T;
  OVERRIDES
    apply1 := Apply1_Lambda_0;
  END;

(* At the point where the lambda appears: *)
VAR lam := NEW(Lambda_0,
    captured_x := a_x,
    b_plus := self.b_plus);
\end{lstlisting}

For free variables that are themselves free in the enclosing
function (i.e., accessed via binding cells), the lambda can share
the same binding cell reference---no special treatment is needed.

\paragraph{Complication: mutable captures.}
If a captured variable is mutated via \texttt{set!} in either the
enclosing function or the lambda, capture-by-value is incorrect.
The standard solution is to ``box'' mutable variables: allocate a
one-element mutable cell and have both the enclosing function and
the lambda access the variable through the cell.  This is the
approach described by Dybvig~[4].  However, \texttt{set!} on
captured variables is rare in practice; the common case (lambdas
passed to \texttt{map}) involves no mutation.  An initial
implementation can reject lambdas that capture mutated variables,
extending the compilability check.

\paragraph{Difficulty: moderate.}
The code generation machinery already exists; the main work is
lambda lifting (generating a fresh type per lambda) and wiring up
the free-variable analysis to distinguish parameter captures from
binding-cell captures.

\subsection{Named \texttt{let}}
\label{sec:gap-named-let}

\paragraph{Impact.}
Seven definitions in \texttt{compiler.scm} use named \texttt{let},
including utility functions like \texttt{iota},
\texttt{unique-symbols}, and \texttt{read-file}.  This is the
second most common gap.

\paragraph{What is needed.}
Named \texttt{let} is syntactic sugar for a self-recursive local
function:
\begin{lstlisting}[language=Scheme]
(let loop ((i 0) (acc '()))
  (if (= i n) acc
      (loop (+ i 1) (cons i acc))))
\end{lstlisting}

\noindent is equivalent to:
\begin{lstlisting}[language=Scheme]
((lambda (loop)
   (set! loop (lambda (i acc)
     (if (= i n) acc
         (loop (+ i 1) (cons i acc)))))
   (loop 0 '()))
 #f)
\end{lstlisting}

\paragraph{Strategy.}
Named \texttt{let} is almost exactly the self-tail-call pattern
that the compiler already handles.  The difference is that the
``function'' is local, not top-level.  The most direct compilation
is to emit a Modula-3 \texttt{LOOP} with the named-let variables
as \texttt{VAR} locals, the initial values as initializers, and the
body as the loop body---exactly the self-tail-call pattern.

Tail calls to the named-let name become parameter reassignment
followed by loop continuation (no \texttt{EXIT}).  Non-tail calls
to the name would require either a nested procedure or a lambda
(and could be rejected by the compilability checker until lambda
support is available).

\begin{lstlisting}[language=Modula3]
(* Named let compiled to LOOP: *)
VAR
  a_i : SchemeObject.T := SchemeLongReal.Zero;
  a_acc : SchemeObject.T := NIL;
BEGIN
  LOOP
    IF (* (= i n) *) THEN
      RETURN a_acc
    ELSE
      VAR t_0, t_1 : SchemeObject.T; BEGIN
        t_0 := SchemeLongReal.Add(a_i, SchemeLongReal.One);
        t_1 := SchemeUtils.Cons(a_i, a_acc, interp);
        a_i := t_0;
        a_acc := t_1;
      END
    END
  END
END
\end{lstlisting}

\paragraph{Difficulty: easy.}
The self-tail-call machinery can be reused almost directly.

\subsection{\texttt{letrec}}
\label{sec:gap-letrec}

\paragraph{Impact.}
Not used in \texttt{compiler.scm}, but common in general Scheme
code for defining mutually recursive local functions.

\paragraph{Strategy.}
\texttt{letrec} binds variables to undefined values, then evaluates
the initializers in the scope of all the bindings (allowing mutual
references).  The compilation is:

\begin{enumerate}
\item Declare all variables as \texttt{VAR} with no initializer
  (defaulting to \texttt{NIL}).
\item Evaluate each initializer and assign to the corresponding
  variable.
\item Execute the body.
\end{enumerate}

This is simpler than \texttt{let} (no shadowing concern, since the
variables are explicitly uninitialized during init evaluation).

\paragraph{Difficulty: easy.}

\subsection{Internal \texttt{define}}
\label{sec:gap-internal-define}

\paragraph{Impact.}
Not used in \texttt{compiler.scm}, but idiomatic in many Scheme
programs for defining local helper functions.

\paragraph{Strategy.}
R5RS specifies that internal defines at the beginning of a body
are equivalent to \texttt{letrec}.  The compiler can desugar
internal defines to \texttt{letrec} during an AST preprocessing
pass, then compile the \texttt{letrec} as described above.

\paragraph{Difficulty: easy,} once \texttt{letrec} is supported.

\subsection{\texttt{cond} with \texttt{=>} Syntax}
\label{sec:gap-cond-arrow}

\paragraph{Impact.}
Not currently used in \texttt{compiler.scm} (the function that
previously used it, \texttt{m3-char}, was eliminated when the
compiler switched to serial-numbered identifiers;
see~\S\ref{sec:impl-idents}).  However, \texttt{cond =>} appears
in general Scheme programs.

\paragraph{What is needed.}
The clause \texttt{(test => proc)} evaluates \texttt{test}; if true,
it calls \texttt{proc} with the test value as argument.  This
requires binding the test result to a temporary and passing it
to the procedure.

\paragraph{Strategy.}
In \texttt{gen-cond}, when a clause has the \texttt{=>} form:

\begin{lstlisting}[language=Modula3]
VAR t_0 : SchemeObject.T; BEGIN
  t_0 := (* evaluate test *);
  IF SchemeBoolean.TruthO(t_0) THEN
    (* call proc with t_0 as argument *)
    RETURN NARROW(proc, SchemeProcedure.T).apply1(interp, t_0)
  ELSE
    (* next clause *)
  END
END
\end{lstlisting}

\paragraph{Difficulty: easy.}
The test value is already computed into a temporary; the only
addition is detecting the \texttt{=>} token and generating a
call instead of evaluating the clause body.

\subsection{\texttt{case}}
\label{sec:gap-case}

\paragraph{Impact.}
Not used in \texttt{compiler.scm}.

\paragraph{Strategy.}
\texttt{case} is a multi-way dispatch on a key value.
Compile as a chain of \texttt{IF}/\texttt{ELSIF} comparisons
using \texttt{eqv?}.  For integer-only keys, a Modula-3
\texttt{CASE} statement could be emitted for better performance.

\paragraph{Difficulty: easy.}

\subsection{\texttt{do}}
\label{sec:gap-do}

\paragraph{Impact.}
Not used in \texttt{compiler.scm}.  Rarely used in MScheme
programs generally.

\paragraph{Strategy.}
\texttt{do} is a general iteration construct.  Compile to a
Modula-3 \texttt{LOOP} with step expressions evaluated into
temporaries and assigned at the end of each iteration, similar
to the self-tail-call pattern.

\paragraph{Difficulty: easy.}

\subsection{Rest Parameters}
\label{sec:gap-rest}

\paragraph{Impact.}
One definition in \texttt{compiler.scm} (\texttt{L}) uses a
rest parameter \texttt{(define (L .\ strs) ...)}.

\paragraph{What is needed.}
A rest parameter collects all arguments beyond the fixed parameters
into a list.  The \texttt{apply} method already receives arguments
as a Scheme list; the compilation needs to split the list into
fixed-parameter locals and a rest-parameter local.

\paragraph{Strategy.}
In the \texttt{Apply\_\emph{foo}} (list-argument) method, after
extracting the fixed parameters with \texttt{SchemeUtils.First}/\texttt{Rest},
the remaining \texttt{rest} variable is bound to the rest parameter.
The \texttt{apply1} and \texttt{apply2} specializations must
construct a list from the excess arguments.

\begin{lstlisting}[language=Modula3]
(* (define (L . strs) body) -- zero fixed params, one rest *)
PROCEDURE Apply_L(self : Compiled_L;
    interp : Scheme.T;
    args : SchemeObject.T) : SchemeObject.T
  RAISES { Scheme.E } =
VAR a_strs := args;  (* entire arg list IS the rest parameter *)
BEGIN
  RETURN ...
END Apply_L;
\end{lstlisting}

\paragraph{Difficulty: easy.}
The main change is in the argument-unpacking code in
\texttt{gen-apply-proc} and \texttt{gen-apply-n-proc}, plus
updating \texttt{proper-list?} to accept and handle dotted
parameter lists.

%----------------------------------------------------------------------
\section{Call Optimization}
\label{sec:gaps-calls}
%----------------------------------------------------------------------

Beyond language coverage, the compiler leaves performance on the
table in how it dispatches calls.  The current call sequence for
a function \texttt{foo} is:

\begin{enumerate}
\item Fetch the callee: \texttt{self.b\_foo.get()} (binding lookup).
\item Type-check: \texttt{NARROW(proc, SchemeProcedure.T)}.
\item Dispatch: \texttt{.apply1(interp, arg)} (virtual method call).
\end{enumerate}

\noindent
Three optimizations can progressively reduce this overhead.

\subsection{Intra-Module Direct Calls}
\label{sec:gap-intra}

This optimization was described in detail in \S\ref{sec:impl-future}.
For calls between functions in the same compilation unit, a pointer
comparison guards a direct call that bypasses both \texttt{NARROW}
and virtual dispatch:

\begin{lstlisting}[language=Modula3]
VAR proc := self.b_foo.get();
BEGIN
  IF proc = self.d_foo THEN
    RETURN Apply1_foo(self.d_foo, interp, arg)
  ELSE
    RETURN NARROW(proc, SchemeProcedure.T).apply1(interp, arg)
  END
END
\end{lstlisting}

The fast path eliminates 2 of the 3 overhead sources (NARROW +
virtual dispatch), leaving only the binding lookup and a well-predicted
branch.  For call-intensive benchmarks like \texttt{fibonacci},
this would approach the performance of hand-written Modula-3.

\subsection{Cross-Module Direct Calls}
\label{sec:gap-cross}

The intra-module optimization requires the caller to hold a typed
reference to the callee's compiled object.  For functions in
\emph{different} compilation units, the compiler does not have such
a reference at code-generation time.

\paragraph{Strategy: registration protocol.}
When module~A calls a function \texttt{bar} defined in module~B,
the fast path requires a \texttt{Compiled\_bar} reference.
This can be provided through a two-phase initialization:

\begin{enumerate}
\item Module~B's \texttt{Install} procedure registers its compiled
  objects in a well-known table (keyed by symbol name).
\item Module~A's \texttt{Install} procedure, called \emph{after}
  module~B's, looks up the compiled object for \texttt{bar} and
  stores a direct reference.
\end{enumerate}

The generated code uses the same guard pattern as intra-module:

\begin{lstlisting}[language=Modula3]
TYPE Compiled_baz = SchemeProcedure.T OBJECT
    b_bar : SchemeEnvironmentBinding.T;
    d_bar : SchemeProcedure.T;  (* untyped -- different module *)
    ...
  END;

(* call site: *)
VAR proc := self.b_bar.get();
BEGIN
  IF proc = self.d_bar THEN
    RETURN self.d_bar.apply1(interp, arg)
  ELSE
    RETURN NARROW(proc, SchemeProcedure.T).apply1(interp, arg)
  END
END
\end{lstlisting}

Note that the cross-module case cannot avoid the virtual dispatch
(since the caller does not know the concrete type), but it can
still skip the \texttt{NARROW} by storing the reference as
\texttt{SchemeProcedure.T} rather than \texttt{SchemeObject.T}.
The pointer comparison still guards against redefinition.

\paragraph{Difficulty: moderate.}
Requires a cross-module registration protocol and careful ordering
of \texttt{Install} calls.

\subsection{General Tail-Call Optimization}
\label{sec:gap-tco}

The compiler currently optimizes only \emph{self}-tail-calls
(direct recursion).  Mutual tail calls (A tail-calls B, B
tail-calls A) and tail calls to arbitrary functions are not
optimized---each becomes a regular call that consumes stack.

\paragraph{Strategy: trampoline.}
The standard technique for general TCO in a language without
goto is the \emph{trampoline}: instead of calling the tail
callee directly, the function returns a thunk (or a tagged value)
that the trampoline loop invokes.  This eliminates stack growth
at the cost of one allocation per tail call.

In the MScheme context, a trampoline could be implemented as a
\texttt{LOOP} in the \texttt{apply} methods that detects a
special ``tail-call request'' return value and re-dispatches:

\begin{lstlisting}[language=Modula3]
(* Tail call to foo becomes: *)
RETURN NEW(TailCall, proc := self.b_foo.get(),
           args := SchemeUtils.List1(arg, interp))
\end{lstlisting}

The trampoline loop in a wrapper would then be:

\begin{lstlisting}[language=Modula3]
VAR result := proc.apply1(interp, arg);
BEGIN
  WHILE ISTYPE(result, TailCall) DO
    WITH tc = NARROW(result, TailCall) DO
      result := NARROW(tc.proc,
          SchemeProcedure.T).apply(interp, tc.args)
    END
  END;
  RETURN result
END
\end{lstlisting}

\paragraph{Difficulty: hard.}
Requires a new return convention and careful integration with
the interpreter's existing call paths.  The allocation overhead
of the \texttt{TailCall} object partially offsets the benefit.
Given that self-tail-call optimization already handles the most
common case (accumulator loops), general TCO is lower priority.

\subsection{Primitive Inlining}
\label{sec:gap-inline}

Calls to known primitives (\texttt{+}, \texttt{-}, \texttt{*},
\texttt{car}, \texttt{cdr}, \texttt{cons}, \texttt{null?}, etc.)
currently go through the full binding-lookup / NARROW / virtual-dispatch
path.  Since primitives are rarely redefined, inlining them would
eliminate per-call overhead entirely.

\paragraph{Strategy.}
Maintain a table of known primitives and their Modula-3
equivalents.  For each known call, emit the Modula-3 code directly:

\begin{lstlisting}[language=Modula3]
(* (+ x y) compiled with primitive inlining: *)
SchemeLongReal.Add(a_x, a_y)

(* vs. current: *)
VAR fn := self.b_plus.get();
    t0, t1 : SchemeObject.T;
BEGIN
  t0 := a_x; t1 := a_y;
  RETURN NARROW(fn, SchemeProcedure.T).apply2(interp, t0, t1)
END
\end{lstlisting}

To preserve redefinition semantics, the inlined code can be guarded
by the same pointer-comparison pattern used for intra-module calls,
with the original primitive object stored as a direct reference.
Alternatively, since primitive redefinition is extremely rare in
practice, the compiler could offer a ``safe mode'' flag
that disables inlining.

\paragraph{Difficulty: moderate.}
Requires enumerating the primitives and their M3 implementations,
handling type dispatch (e.g., \texttt{+} on integers vs.\ reals),
and deciding on the redefinition guard policy.

%----------------------------------------------------------------------
\section{Self-Compilation Summary}
\label{sec:gaps-selfcomp}
%----------------------------------------------------------------------

The following table summarizes the self-compilation gap analysis.
The compiler's own source (\texttt{compiler.scm}) serves as a
useful stress test.  The distribution of blocking constructs
mirrors that of the CSP compiler (\S\ref{sec:impl-cspc}):
\texttt{lambda} dominates, followed by named \texttt{let}.

\begin{center}
\begin{tabular}{llc}
\toprule
\textbf{Blocking construct} & \textbf{Typical functions} & \textbf{Approx.} \\
\midrule
\texttt{lambda} (in \texttt{map}/\texttt{filter} args)
  & \texttt{gen-let}, \texttt{gen-call}, \texttt{gen-type-decl},
    \texttt{gen-apply-proc}, &  \\
  & \texttt{gen-install-proc}, \texttt{gen-m3},
    \texttt{compile-file}, \ldots & 12+ \\
\texttt{lambda} + named \texttt{let}
  & \texttt{free-variables}, \texttt{unique-symbols},
    \texttt{unique-constants}, & \\
  & \texttt{has-self-tail-call?}, \texttt{constant-index} & 5 \\
Named \texttt{let} only
  & \texttt{iota}, \texttt{read-file},
    \texttt{build-indexed-map} & 3 \\
Rest parameter
  & \texttt{L} & 1 \\
\midrule
\textbf{Total skipped} & & $\sim$\textbf{21} \\
\bottomrule
\end{tabular}
\end{center}

Adding \texttt{lambda} support would unblock $\sim$12 definitions.
Adding both \texttt{lambda} and named \texttt{let} would unblock
$\sim$20, leaving only \texttt{L} (rest parameter)---an easy
addition that would complete full self-compilation.  The remaining
gaps (\texttt{letrec}, \texttt{case}, \texttt{do}, internal
\texttt{define}) are not used in the compiler itself but appear
in general Scheme programs; their frequency in the CSP compiler
stress test is documented in \S\ref{sec:impl-cspc}.

%----------------------------------------------------------------------
\section{Recommended Priorities}
\label{sec:gaps-priorities}
%----------------------------------------------------------------------

Based on the impact analysis above, the recommended implementation
order is:

\begin{enumerate}
\item \textbf{Named \texttt{let}} --- easy, reuses existing
  self-tail-call machinery, unblocks 2 definitions immediately and
  5 more when combined with lambda.

\item \textbf{\texttt{lambda}} --- moderate difficulty, unblocks 13+
  definitions, required for self-compilation.  Start with the
  common case (lambdas that capture no mutable variables) and
  extend later.

\item \textbf{Intra-module direct calls} --- moderate, gives the
  largest performance improvement for compiled code, especially
  for recursive benchmarks.

\item \textbf{Rest parameters} --- easy, one-off addition,
  completes self-compilation.

\item \textbf{\texttt{cond =>}} --- easy, one-off addition.

\item \textbf{\texttt{letrec} and internal \texttt{define}} ---
  easy, low urgency.

\item \textbf{Cross-module direct calls} --- moderate, deferred
  until multiple compiled modules are common.

\item \textbf{Primitive inlining} --- moderate, largest potential
  speedup but requires careful design.

\item \textbf{General tail-call optimization} --- hard, lowest
  priority since self-tail-calls cover the common case.
\end{enumerate}

Steps 1--2 would achieve full self-compilation of the compiler.
Steps 1--4 would close all gaps observed in real MScheme programs
(cf.\ the CSP compiler stress test, \S\ref{sec:impl-cspc}).
Steps 5--9 are optimizations that become relevant as the compiled
code base grows.

%======================================================================
\chapter*{References}
\addcontentsline{toc}{chapter}{References}
%======================================================================

\begin{enumerate}
\item The companion report: \emph{MScheme Internals: A Technical
  Report on the Interpreter Implementation}, March 2026.  Describes
  the binding optimization, evaluator loop, environment system, and
  primitive dispatch in full detail.

\item Peter Norvig, ``JScheme: Scheme in Java,''
  \url{https://norvig.com/jscheme.html}, 1998.
  The origin of the MScheme design.

\item Marc Feeley and Guy Lapalme, ``Using Closures for Code
  Generation,'' \emph{Computer Languages}, 12(1):47--66, 1987.
  Describes the technique of compiling Scheme to closures, which is
  closely related to the approach proposed here.

\item R.\ Kent Dybvig, \emph{Three Implementation Models for Scheme},
  PhD thesis, University of North Carolina at Chapel Hill, 1987.
  Describes heap-based, stack-based, and string-based implementation
  models for Scheme, including closure conversion and environment
  representation strategies relevant to this design.

\item\label{ref:nelson} Greg Nelson (ed.), \emph{Systems Programming with Modula-3},
  Prentice Hall, 1991.

\item Samuel P.\ Harbison, \emph{Modula-3}, Prentice Hall, 1992.

\item Richard Kelsey, William Clinger, and Jonathan Rees (eds.),
  ``Revised$^5$ Report on the Algorithmic Language Scheme,'' 1998.
  \url{https://www.schemers.org/Documents/Standards/R5RS/}.
\end{enumerate}

\end{document}
